% LaTeX companion of 02-random-variables.typ. Both formats are editable.
\chapter{random variable}\label{random-variable}

\section{\texorpdfstring{random variable and generated \(\mathit{\sigma}\)-algebra}{random variable and generated \textbackslash mathit\{\textbackslash sigma\}-algebra}}\label{random-variable-and-generated-sigma-algebra}

\subsection{random variable: 即 prob space 上的一个 Borel measurable function}\label{random-variable-ux5373-prob-space-ux4e0aux7684ux4e00ux4e2a-borel-measurable-function}

\begin{definition}{\kn{random variable}}

对于 \knref{probability space} \((\Omega,\mathcal{F},\mathbb{P})\), 一个 random variable 是一个 \textbf{\knref{$(\mathcal{M},\mathcal{N})$-measurable function}（Borel measurable function）} \(\mathit{X}:\Omega\rightarrow\mathbb{R}\).

\end{definition}

\begin{qlremark}{Remark}

回顾 (Borel) measurable function 的定义: 即对于任意 \(\mathit{B} \in \mathcal{B}(\mathbb{R})\), 有 \(\mathit{X}^{- 1}(\mathit{B}) \in \mathcal{F}\); 由于是在 real space 上, Borel measurable 也可以简化为: 对于任意的 open ray \(( - \infty,\mathit{a})\), 有 \(\mathit{X}^{- 1}(( - \infty,\mathit{a})) \in \mathcal{F}\).

例如: \(\mathit{X}(\mathit{\omega}) = \mathit{\omega}^{2}\) 是一个 (Borel) measurable function, 因为对于任意 open ray \(( - \infty,\mathit{a})\), \(\mathit{X}^{- 1}(( - \infty,\mathit{a}))\) 是一个 interval (either \(( - \sqrt{\mathit{a}},\sqrt{\mathit{a}})\) or \(\varnothing\)).

并且 recall: 任意的 continuous function 都是 (Borel) measurable function. (所以之后我们看到 continuous function 就知道它一定是一个 random variable).

\end{qlremark}

\begin{qlremark}{Remark}

Random variable (Borel measurable function) 本质上是对 样本空间中的事件进行 weights assignment.

原本, 一个 prob space 只能表示: \textbf{有多少个 outcomes, 每个 outcome 的概率是多少} (对于无限的 sample space, 一些 outcomes 组合, 即事件, 的概率是多少).

而 random variable 就是\textbf{给这些 outcomes 以 weights}: 比如

\begin{example}

在一副扑克牌中抽一张, 那么每一张的概率都是 \(1/52\). 但是如果我想知道抽到的牌的价值, 那么我就需要一个 random variable \(\mathit{X}\) 来给每一张牌赋予一个 weight; 比如说, \(\mathit{X}(\text{A}) = 1\), \(\mathit{X}(2) = 2\), \(\cdots\), \(\mathit{X}(\text{K}) = 13\).\\

\end{example}

\textbf{Question}: 为什么对于 assigning weights 的行为, 我们一定要求它是一个 (Borel) measurable function 呢? 普通的给 outcomes 以 weights 的行为 (任意一个 function \(\mathit{f}:\Omega\rightarrow\mathbb{R}\)) 不行吗?

\textbf{Answer}: 因为我们想要知道 \(\mathit{X}\) 的值在一个范围 \(\mathit{B}\), 比如 \(\lbrack 1,6\rbrack\) 之内的 probability (which is defined on 原本的 prob measure \(\mathbb{P}\) 上的) 是多少. 而只有每一个这样的范围集合的 preimage \(\mathit{X}^{- 1}(\mathit{B})\) 都是 \(\mathcal{F}\) 中的事件, 我们才能够通过 \(\mathbb{P}\) 来计算这个 probability; that is, 我们需要这个 \(\mathit{X}\) 是一个 measurable function.

因为这一切都是建立在我们以 measure theory 的 framework 来定义 probability 的基础上的.\\

\end{qlremark}

下面是一个简单的 proposition: 我们可以把多个 random variables 以一种 Borel measurable 的方式组合起来, 那么就成为一个新的 random variable.

\begin{proposition}{\kn{以 Borel measurable 的方式组合起多个 random variables}}

Prob space \((\Omega,\mathcal{F},\mathbb{P})\) 上的 random variables \(\mathit{X}_{1},\mathit{X}_{2},\ldots,\mathit{X}_{\mathit{k}}:\Omega\rightarrow\mathbb{R}\), 任取 Borel measurable function \(\mathit{g}:\mathbb{R}^{\mathit{k}}\rightarrow\mathbb{R}\),

\[\begin{matrix}
{\mathit{X}:\Omega} & {\rightarrow\mathbb{R}} \\
\mathit{\omega} & {\mapsto\mathit{g}(\mathit{X}_{1}(\mathit{\omega}),\mathit{X}_{2}(\mathit{\omega}),\ldots,\mathit{X}_{\mathit{k}}(\mathit{\omega}))}
\end{matrix}\]

也是一个 random variable.

\end{proposition}

\begin{proof}

我们在 measure theory 中证明过: 对于任意的 finite seq of Borel measureable functions \((\mathit{f}_{\mathit{i}}:\Omega\rightarrow\mathbb{R})_{\mathit{i} = 1}^{\mathit{k}}\), 其各作为一个维度组成的函数 \(\mathit{f} = (\mathit{f}_{1},\cdots,\mathit{f}_{\mathit{k}})\) 也是一个 Borel measurable function (from \(\Omega\) 到 \(\mathbb{R}^{\mathit{k}}\)).\\
而这里的 \(\mathit{X}\) 就是 \(\mathit{g} \circ \mathit{f}\), 是两个 Borel measurable function 的 composition, 因而也是一个 Borel measurable function (即 random variable).

\end{proof}

\begin{qlremark}{Remark}

以一种 Borel measurable 的方式组合多个 random variables, 这个函数也仍然是一个 random variable.

这个 proposition 看似没有什么卵用, 但是它其实是我们之后构造 random variable 的一个重要工具. 很多我们默认合法的行为, 它都作为基础.

比如我们:

\begin{itemize}
\item
  定义了两个 random variable \(\mathit{X}\) 和 \(\mathit{Y}\), 那么它们的和 \(\mathit{X} + \mathit{Y}\) 以及积 \(\mathit{X} \cdot \mathit{Y}\) 也是 random variables.
\item
  更加复杂的: 标量变换 \(\mathit{g}(\mathit{X}) = \sin(\mathit{X}),h(\mathit{X}) = \mathit{e}^{\mathit{X}}\), 向量变换 \(\mathit{g}(\mathit{X}_{1},\mathit{X}_{2}) = \mathit{X}_{1}^{2} + \mathit{X}_{2}^{2}\), 等等, 都是 random variables.
\item
  统计学中, 我们有一堆样本, \textbf{每一个都是随机采样得到的 (即一个 random variable \(\mathit{X}_{\mathit{i}}\))}. 因而在观测之前, 我们 把它看作一个 random variable \(\mathit{X}_{\mathit{i}}\), 然后在观测之后取到它的一个实际值 \(\mathit{x}_{\mathit{i}}\).

  通常我们假设所有样本都是 \textbf{independent and identically distributed (i.i.d., 独立同分布)} 的, 即所有的 \(\mathit{X}_{1},\mathit{X}_{2},\ldots,\mathit{X}_{\mathit{n}}\) 都是 independent 的, 并且服从同一个 distribution (discuss later\ldots).

  然后, 我们就可以从这些样本中计算各种统计量, 比如平均值 \(\bar{\mathit{X}} := \frac{1}{\mathit{n}}\sum_{\mathit{i} = 1}^{\mathit{n}}\mathit{X}_{\mathit{i}}\), 再比如 order statistic, sample variance 等等, 从而通过大数定律等工具的辅佐, 来估计这个 distribution 到底是什么样的一个 distribution.

  这里这些统计量, 由于是以 Borel measurable 的方式组合起多个 random variables, 也是 random variables. 所以这个 proposition 其实也奠定了统计学中各种统计量的合法性.
\end{itemize}

\end{qlremark}

\subsection{\texorpdfstring{由一个 random variable generate 的 \(\mathit{\sigma}\)-algebra}{由一个 random variable generate 的 \textbackslash mathit\{\textbackslash sigma\}-algebra}}\label{ux7531ux4e00ux4e2a-random-variable-generate-ux7684-mathitsigma-algebra}

\begin{definition}{\kn{$\mathit{\sigma}$-algebra generated by a random variable (measurable
function)}}

对于 random variable \(\mathit{X}:\Omega\rightarrow\mathbb{R}\), 其 generate 的 \(\mathit{\sigma}\)-algebra 定义为

\[\mathit{\sigma}(\mathit{X}) := \left\{ {\mathit{X}^{- 1}(\mathit{B}):\mathit{B} \in \mathcal{B}(\mathbb{R})} \right\} \subseteq \mathcal{F}\]

\end{definition}

这个集合 \(\mathit{\sigma}(\mathit{X})\) 是一个 \(\mathit{\sigma}\)-algebra, 是一个 trivial truth. 因为 里面所有的 set 都是 \(\mathit{X}\) 的 preimage, 而 \(\mathit{X}\) 是一个 measurable function, 因而任意的 \(\mathit{X}^{- 1}(\mathit{B}) \in \mathcal{F}\).

{\emergencystretch=2em
这个 \(\mathit{\sigma}(\mathit{X})\) 也是让 \(\mathit{X}\) 成为一个 measurable function (即 random variable) 的最小的 \(\mathit{\sigma}\)-algebra. 它的意义是: \textbf{当我们仅仅观察到 \(\mathit{X}\) 的取值时, 我们所掌握的关于 underlying sample space \(\Omega\) 的全部 "information".}\par}


\begin{qlremark}{Remark}

我们知道, \textbf{一个 \(\mathit{\sigma}\)-algebra \(\mathcal{F}\) 可以看作是一个 information structure}: 它是对 underlying sample space \(\Omega\) 进行划分的一个工具.\\
最大的 \(\mathit{\sigma}\)-algebra \(\mathcal{P}(\Omega)\), 即 power set, 是 最细的划分, 也是信息量最大的 \(\mathit{\sigma}\)-algebra. 它把每一个元素都区分开来. 对于 discrete 的 sample space, 通常我们的 probability measure 就是定义在 power set 上的, 正是因为我们 希望区分每一个 outcome.\\
而给定一个 prob space \((\Omega,\mathcal{F},\mathbb{P})\), 原本的 prob measure 所 base on 的 \(\mathit{\sigma}\)-algebra \(\mathcal{F}\) 就成为了一个基准的 information 的颗粒度: 哪些事件是我们可以区分, 即可以在上面定义 probability 的.

而一个 random variable \(\mathit{X}\) 可能会把多个 outcomes 映射到同一个值上, 也即 把多个 events 映射到同一个值上, 从而把这些 events 区分不开来. 因而, \textbf{random variable \(\mathit{X}\) generate 的 \(\mathit{\sigma}\)-algebra \(\mathit{\sigma}(\mathit{X})\) 就可能是一个更粗的划分}. 当我们取这些 events 所映射到的同一个值的 preimage 的时候, 取到的是 它们的 union, 也就是它们中最大的一个.

比如

\begin{example}

掷一次骰子, 所有可能结果为 sample space:

\[\Omega = \left\{ {1,2,3,4,5,6} \right\}\]

而我们定义一个 random variable \(\mathit{X}\) 来表示掷出的点数的奇偶性, 即:

\[\mathit{X}(\mathit{\omega}) = \left\{ \begin{matrix}
{1,} & {\mathit{\omega} \in \left\{ {2,4,6} \right\}} \\
{0,} & {\mathit{\omega} \in \left\{ {1,3,5} \right\}}
\end{matrix} \right.\]

那么它 generate 的 \(\mathit{\sigma}\)-algebra:

\[\mathit{\sigma}(\mathit{X}) = \left\{ {\varnothing,\left\{ {1,3,5} \right\},\left\{ {2,4,6} \right\},\Omega} \right\}\]

因为 \(1,3,5\) 这些事件都被映射到了同一个值 \(0\) 上, 即 \(\left\{ 1 \right\},\left\{ 3 \right\},\left\{ 5 \right\},\left\{ {1,3} \right\}\) 这些事件都被 \(\left\{ {1,3,5} \right\}\) 这个 事件吸收了, 它们在这个 random variable \(\mathit{X}\) 的映射下是区分不开的. 这就代表了 这个 random variable 所蕴含的信息颗粒度: \textbf{我们只能区分掷出的点数是奇数还是偶数, 但是无法区分具体的点数是多少.}\\
至此我们知道: \textbf{一个 random variable \(\mathit{X}\) 的 generate 的 \(\mathit{\sigma}\)-algebra \(\mathit{\sigma}(\mathit{X})\) 表示了这个 random variable 对 underlying sample space \(\Omega\) 所能区分的 information granularity (信息颗粒度)}.\\

\end{example}

\end{qlremark}

除了能表示 information granularity 以外, \(\mathit{\sigma}(\mathit{X})\) 还有什么实际用处呢:

\begin{itemize}
\item
  之后讨论关于 independence of two random variables \(\mathit{X}\) 和 \(\mathit{Y}\) 时, 可以用它来作严格定义, 并且展现 independence 的本质: 两个 RV 的 independence 实际表示它们蕴含的信息颗粒度之间没有任何 overlap
\item
  用以定义 conditional expectation: \(\left. \mathbb{E}\lbrack\mathit{Y} \middle| \mathit{\sigma}(\mathit{X})\rbrack \right.\), 通常 简写为 \(\left. \mathbb{E}\lbrack\mathit{Y} \middle| \mathit{X}\rbrack \right.\), 但是其实 \(\left. \mathbb{E}\lbrack\mathit{Y} \middle| \mathit{\sigma}(\mathit{X})\rbrack \right.\) 是一个更加直观的事情, 表述了一种 Partial Averaging 的概念.
\item
  Stochastic Process 中有众多应用. 比如 stopping time 的定义.
\end{itemize}

至此关于 random variable, 我们已经讨论了很多 general 的 刻画. 现在我们来讲一点实用的:

\section{distributions}\label{distributions}

\subsection{distribution and cumulative distribution function}\label{distribution-and-cumulative-distribution-function}

\begin{definition}{\kn{probability distribution}}

对于 random variable \(\mathit{X}:\Omega\rightarrow\mathbb{R}\) on \(\Omega,\mathcal{F},\mathbb{P}\), 其 probability distribution 是 \(\mathit{X}\) 对于 \(\mathbb{P}\) 的 pushforward measure, 记作 \(\mathbb{P}^{\mathit{X}}:\mathcal{B}(\mathbb{R})\rightarrow\lbrack 0,1\rbrack\). 即

\[\mathbb{P}^{\mathit{X}}(\mathit{B}) := \mathbb{P}(\mathit{X}^{- 1}(\mathit{B})),\quad\forall\mathit{B} \in \mathcal{B}(\mathbb{R})\]

我们 write:

\[\mathit{X} \sim \mathbb{P}^{\mathit{X}}\]

\end{definition}

\begin{qlremark}{Remark}

看着有点不直观. 实则我们拆解:

\begin{itemize}
\item
  \(\mathit{X}^{- 1}(\mathit{B})\) 即: 有多少 samples (即哪一个事件 \(\mathit{E} \in \mathcal{F}\)) 在这个 RV \(\mathit{X}\) 的映射下落入 \(\mathit{B}\) 这个区间?
\item
  \(\mathbb{P}(\mathit{X}^{- 1}(\mathit{B}))\) 即: 这个事件 \(\mathit{E}\) 的概率是多少?
\end{itemize}

接下来的定义和例子会更加直观.

\end{qlremark}

\begin{definition}{\kn{distribution function (也称 \textbf{cumulative distribution function,
cdf})}}

对于 random variable \(\mathit{X}:\Omega\rightarrow\mathbb{R}\), 其 \textbf{cumulative distribution function} 是 \(\mathbb{P}^{\mathit{X}}\) 这一函数的 distribution function, 记作 \(\mathit{F}_{\mathit{X}}\). 即对于任意 \(\mathit{x} \in \mathbb{R}\), 有

\[\mathit{F}_{\mathit{X}}(\mathit{x}) := \mathbb{P}^{\mathit{X}}(( - \infty,\mathit{x}\rbrack) = \mathbb{P}(\mathit{X}^{- 1}(( - \infty,\mathit{x}\rbrack))\]

\end{definition}

\begin{qlremark}{Remark}

\(\mathit{F}_{\mathit{X}}(\mathit{x})\) 即 \(\mathbb{P}(\left\{ {\mathit{X} \leq \mathit{x}} \right\})\) 我们也会简写为 \(\mathbb{P}(\mathit{X} \leq \mathit{x})\).\\
\(\mathit{F}_{\mathit{X}}(\mathit{x})\) 的意义是: random variable \(\mathit{X}\) 的值不超过 \(\mathit{x}\) 的概率.\\

\end{qlremark}

\begin{example}[(geometric probability)]

Fix \(\mathit{a},\mathit{b} > 0\). 我们在正方形区域 \(\Omega = \lbrack 0,\mathit{a}\rbrack \times \lbrack 0,\mathit{b}\rbrack\) 上均匀地随机选取一个点 \((\mathit{x},\mathit{y})\).\\
我们 define 随机变量: \(\mathit{X}:\Omega\rightarrow\mathbb{R}\) by \(\mathit{X}(\mathit{x},\mathit{y}) = \mathit{x}\). 求 \(\mathit{X}\) 的 distribution function \(\mathit{F}_{\mathit{X}}\).\\
这很简单: 对于 \(\mathit{x} \in \lbrack 0,\mathit{a}\rbrack\),

\[\mathit{F}_{\mathit{X}}(\mathit{x}) = \mathbb{P}(\mathit{X} \leq \mathit{x}) = \frac{\mathit{m}(\lbrack 0,\mathit{x}\rbrack \times \lbrack 0,\mathit{b}\rbrack)}{\mathit{m}(\lbrack 0,\mathit{a}\rbrack \times \lbrack 0,\mathit{b}\rbrack)} = \frac{\mathit{x} \cdot \mathit{b}}{\mathit{a} \cdot \mathit{b}} = \frac{\mathit{x}}{\mathit{a}},\quad 0 \leq \mathit{x} \leq \mathit{a}\]

因为

\[\mathit{F}_{\mathit{X}}(\mathit{x}) = \left\{ \begin{matrix}
{0,} & {\mathit{x} < 0} \\
{\frac{\mathit{x}}{\mathit{a}},} & {0 \leq \mathit{x} \leq \mathit{a}} \\
{1,} & {\mathit{x} > \mathit{a}}
\end{matrix} \right.\]

注意: 在这个例子中, probability measure \(\mathit{P}\) 是 Lebesgue measure on \(\lbrack 0,\mathit{a}\rbrack \times \lbrack 0,\mathit{b}\rbrack\). 很多时候我们在 计算 RV 的 distribution function 的时候, 都是直接直观地计算 \(\mathit{P}(\mathit{X} \leq \mathit{x})\). 但是在稍微复杂一些 的例子中, 我们需要对各个条件的 formalization 更清楚一点.

\end{example}

% qlnotes: kind=theorem; id=thm-02-random-variables-distribution-function
\begin{theorem}{\kn{distribution function 的性质}}\label{thm-02-random-variables-distribution-function}

对于 random variable \(\mathit{X}:\Omega\rightarrow\mathbb{R}\), 其 distribution function \(\mathit{F}_{\mathit{X}}\) 满足:

\begin{itemize}
\item
  \(\mathit{F}_{\mathit{X}}\) 是 non-decreasing (non-strictly increasing) 的.
\item
  \(\lim_{\mathit{x}\rightarrow + \infty}\mathit{F}_{\mathit{X}}(\mathit{x}) = 1\), \(\lim_{\mathit{x}\rightarrow - \infty}\mathit{F}_{\mathit{X}}(\mathit{x}) = 0\).
\item
  \(\mathit{F}_{\mathit{X}}\) 是 right-continuous 的. 即对于任意 \(\mathit{x}_{0}\), 有 \(\lim_{\mathit{x}\rightarrow\mathit{x}_{0}^{+}}\mathit{F}_{\mathit{X}}(\mathit{x}) = \mathit{F}_{\mathit{X}}(\mathit{x}_{0})\).
\item
  对于任意 \(\mathit{x}_{0}\), 有 \(\lim_{\mathit{x}\rightarrow\mathit{x}_{0}^{-}}\mathit{F}_{\mathit{X}}(\mathit{x}) = \mathit{P}(\mathit{X} < \mathit{x}_{0})\). 并且 \(\mathbb{P}(\mathit{X} = \mathit{x}_{0}) = \lim_{\mathit{x}\rightarrow\mathit{x}_{0}^{+}}\mathit{F}_{\mathit{X}}(\mathit{x}) - \lim_{\mathit{x}\rightarrow\mathit{x}_{0}^{-}}\mathit{F}_{\mathit{X}}(\mathit{x})\).
\item
  对于任意 \(\mathit{x}_{1} < \mathit{x}_{2}\), 有 \(\mathit{P}(\mathit{X} \in (\mathit{x}_{1},\mathit{x}_{2}\rbrack) = \mathit{F}_{\mathit{X}}(\mathit{x}_{2}) - \mathit{F}_{\mathit{X}}(\mathit{x}_{1})\).
\end{itemize}

\end{theorem}

\begin{proof}

其他都显然. right-continuity 是源自 measure 的 continuity from above, 我们证明一下: 考虑一个单调递减序列 \(\left\{ \mathit{x}_{\mathit{n}} \right\} \downarrow \mathit{x}_{0}\), 令 seq of events \(\mathit{A}_{\mathit{n}} = \left\{ {\mathit{X} \leq \mathit{x}_{\mathit{n}}} \right\}\), 注意这是一个嵌套递减的 set seq. 通过 \(\mathbb{R}\) 的 completeness 容易证明:

\[\mathit{A}_{0} = \bigcap\limits_{\mathit{n} = 1}^{\infty}\mathit{A}_{\mathit{n}} = \left\{ {\mathit{\omega}:\mathit{X}(\mathit{\omega}) \leq \mathit{x}_{0}} \right\}\]

由 measure 的 continuity from above, \(\mathbb{P}(\mathit{A}_{0}) = \lim_{\mathit{n}\rightarrow\infty}\mathbb{P}(\mathit{A}_{\mathit{n}})\). 也即 \(\lim_{\mathit{x}\rightarrow\mathit{x}_{0}^{+}}\mathit{F}_{\mathit{X}}(\mathit{x}) = \mathbb{P}(\mathit{X} \leq \mathit{x}_{0})\).\\
而下面一条 \(\lim_{\mathit{x}\rightarrow\mathit{x}_{0}^{-}}\mathit{F}_{\mathit{X}}(\mathit{x}) = \mathbb{P}(\mathit{X} < \mathit{x}_{0})\) 同理是源自 measure 的 continuity from below.\\
那么 \(\mathbb{P}(\mathit{X} = \mathit{x}_{0}) = \mathbb{P}(\mathit{X} \leq \mathit{x}_{0}) - \mathbb{P}(\mathit{X} < \mathit{x}_{0})\) 是 natural 的 (by def).

\end{proof}

\begin{qlremark}{Remark}

我们此时心里会默认一件事情: 一旦知道了 RV \(\mathit{X}\) 的 distribution function \(\mathit{F}_{\mathit{X}}\), 就知道了 \(\mathit{X}\) 的完整 distribution \(\mathit{P}^{\mathit{X}}\). 这个直觉是对的.\\
回忆一下: 我们在 measure theory 中, 证明 Carathéodory extension theorem 的时候, 证明过一个 \(\mathit{\pi} - \mathit{\lambda}\) theorem, 以它作为工具才 把 premeasure 从一个 algebra extend 到一个 measure on the generated \(\mathit{\sigma}\)-algebra:

\begin{theorem}{\kn{Dynkin's $\mathit{\pi} - \mathit{\lambda}$ theorem}}

设 \(\mathcal{P}\) 是一个 \(\mathit{\pi}\)-system (closed under finite intersection), \(\mathcal{L}\) 是一个 \(\mathit{\lambda}\)-system (closed under complementation and countable disjoint union), 如果 \(\mathcal{P} \subseteq \mathcal{L}\), 那么 \(\mathit{\sigma}(\mathcal{P}) \subseteq \mathcal{L}\).

\end{theorem}

我们考虑所有的 half-open interval \(( - \infty,\mathit{x}\rbrack\) 组成的集合 \(\mathcal{G}\), 这个集合 generate the Borel \(\mathit{\sigma}\)-algebra, 并且 还是一个 \(\mathit{\pi}\)-system, 即: 它 closed under finite intersection:

\[( - \infty,\mathit{a}\rbrack \cap ( - \infty,\mathit{b}\rbrack \in \mathcal{G},\quad\forall\mathit{a},\mathit{b} \in \mathbb{R}\]

然后考虑

\[\mathcal{L} := \left\{ {\mathit{E} \in \mathit{\sigma}(\mathcal{G}):\mathbb{P}_{1}(\mathit{E}) = \mathbb{P}_{2}(\mathit{E})} \right\}\]

即两个 measure 一致的 collection of sets. 容易证明, 这个 \(\mathcal{L}\) 是一个 \(\mathit{\lambda}\)-system, 并且 \(\mathcal{G} \subset \mathcal{L}\).\\
从而,

\[\mathbb{B}(\mathbb{R}) = \mathit{\sigma}(\mathcal{G}) \subseteq \mathcal{L}\]

并且 \(\mathbb{P}_{1}(\mathit{E}) = \mathbb{P}_{2}(\mathit{E})\) 对于任意 \(\mathit{E} \in \mathbb{B}(\mathbb{R})\).

因而:

\begin{corollary}{\kn{distribution function $\mathit{F}_{\mathit{X}}$ determines
distribution $\mathbb{P}^{\mathit{X}}$}}

对于 random variable \(\mathit{X}:\Omega\rightarrow\mathbb{R}\), 其 distribution function \(\mathit{F}_{\mathit{X}}\) 确定了 distribution \(\mathbb{P}^{\mathit{X}}\); 即 如果两个 random variable \(\mathit{X}\) 和 \(\mathit{Y}\) 满足 \(\mathit{F}_{\mathit{X}} = \mathit{F}_{\mathit{Y}}\), 那么 \(\mathbb{P}^{\mathit{X}} = \mathbb{P}^{\mathit{Y}}\).

\end{corollary}

\end{qlremark}

\subsection{discrete random variable 与 probability mass function}\label{discrete-random-variable-ux4e0e-probability-mass-function}

我们基本上研究两类 random variable: discrete random variable 和 continuous random variable.

\begin{definition}{\kn{discrete random variable}}

对于 random variable \(\mathit{X}:\Omega\rightarrow\mathbb{R}\), 如果存在一个 countable set \(\mathit{S} \subseteq \mathbb{R}\) 使得 \(\mathbb{P}(\mathit{X} \in \mathit{S}) = 1\), 那么 \(\mathit{X}\) 是一个 discrete random variable.

\end{definition}

\begin{qlremark}{Remark}

就是说, \textbf{\(\mathit{X}\) 的 range a.s. 是 countable 的} (允许 null set 上 的值不在这个 countable set 里, 但是这是 measure theory 下的概念, 所以这些都可以忽略不计).

\end{qlremark}

\begin{qlremark}{Remark}

对于 discrete random variable \(\mathit{X}\), 其 distribution function \(\mathit{F}_{\mathit{X}}\) 是一个 step function, 只有 countable 多个 jump discontinuity, 每个 jump 的大小就是 \(\mathit{X}\) 在那个点的 probability mass.\\
因而, 对于 discrete random variable \(\mathit{X}\), 通常考虑其 \textbf{probability mass function (pmf)} \(\mathit{p}_{\mathit{X}}\) 更加有用:

\begin{definition}{\kn{probability mass function (pmf)}}

对于 discrete random variable \(\mathit{X}\), 其 pmf 定义为

\[\mathit{p}_{\mathit{X}}(\mathit{x}) := \mathbb{P}(\mathit{X} = \mathit{x}),\quad\forall\mathit{x} \in \mathbb{R}\]

\end{definition}

这个 pmf 就是 \(\mathit{F}_{\mathit{X}}\) 的 jump size function. 当然, 根据定义明显可得:

\[\sum\limits_{\mathit{x}}\mathit{p}_{\mathit{X}}(\mathit{x}) = 1\]

\end{qlremark}

值得一提的是, 这个 \textbf{pmf 其实就是 \(\mathit{X}\) 的 distribution \(\mathbb{P}^{\mathit{X}}\) 对于 counting measure (for range of \(\mathit{X}\))的 Radon-Nikodym derivative:}

\begin{proposition}{\kn{pmf 是 distribution 对于 counting measure $\mathit{\mu}_{\mathit{S}}$
的 Radon-Nikodym 导数}}

对于一个 discrete random variable \(\mathit{X}\), 假设其 range 是 countable set \(\mathit{S}\), 则我们定义 counting measure \(\mathit{\mu}_{\mathit{S}}\) on \(\mathit{S}\) 为:

\[\mathit{\mu}_{\mathit{S}}(\mathit{A}) := \sum\limits_{\mathit{x}_{\mathit{i}} \in \mathit{S}}\mathbf{1}_{\mathit{A}}(\mathit{x}_{\mathit{i}})\]

那么, distribution \(\mathbb{P}^{\mathit{X}} \ll \mathit{\mu}_{\mathit{S}}\) (绝对连续), 并且有:

\[\mathit{p}_{\mathit{X}} = \frac{\mathit{d}\mathbb{P}^{\mathit{X}}}{\mathit{d}\mathit{\mu}_{\mathit{S}}}\]

\end{proposition}

\begin{proof}

这很容易证明. 首先, 这个 counting measure 即: \(\mathit{A}\) 中有多少个点在 \(\mathit{S}\) 中.\\
那么显然, 如果 \(\mathit{\mu}_{\mathit{S}}(\mathit{A}) = 0\), 即 \(\mathit{A}\) 中没有点在 \(\mathit{S}\) 中, 那么 \(\mathbb{P}^{\mathit{X}}(\mathit{A}) = 0\). 因而 \(\mathbb{P}^{\mathit{X}} \ll \mathit{\mu}_{\mathit{S}}\).\\
且对于任意 Borel set \(\mathit{A}\),

\[\mathbb{P}^{\mathit{X}}(\mathit{A}) = \sum\limits_{\mathit{x}_{\mathit{i}} \in \mathit{A}}\mathbb{P}(\mathit{X} = \mathit{x}_{\mathit{i}}) = \sum\limits_{\mathit{x}_{\mathit{i}} \in \mathit{A} \cap \mathit{S}}\mathit{p}_{\mathit{X}}(\mathit{x}_{\mathit{i}}) = \int_{\mathit{A}}\mathit{p}_{\mathit{X}}(\mathit{x})\mathit{d}\mathit{\mu}_{\mathit{S}}(\mathit{x})\]

\end{proof}

\subsection{continuous random variable 与 probability density function}\label{continuous-random-variable-ux4e0e-probability-density-function}

相对应 discrete random variable, 我们定义 continuous random variable:

\begin{definition}{\kn{continuous random variable}}

我们称一个 random variable \(\mathit{X}:\Omega\rightarrow\mathbb{R}\) 是一个 \textbf{continuous random variable}, 如果它的 cdf \(\mathit{F}_{\mathit{X}}\) 是一个 \textbf{absolutely continuous function.}

\end{definition}

\begin{qlremark}{Remark}

回顾 absolutely continuous function 的定义, 这是一个比较复杂的定义: 对于任意 \(\mathit{\varepsilon} > 0\), 都存在 \(\mathit{\delta} > 0\) 使得任取 finite collection of disjoint intervals \(\left\{ {(\mathit{a}_{\mathit{i}},\mathit{b}_{\mathit{i}})} \right\}\), 都满足:

\[\left. \sum\limits_{\mathit{i}}(\mathit{b}_{\mathit{i}} - \mathit{a}_{\mathit{i}}) < \mathit{\delta}\Longrightarrow\sum\limits_{\mathit{i}} \middle| \mathit{F}_{\mathit{X}}(\mathit{b}_{\mathit{i}}) - \mathit{F}_{\mathit{X}}(\mathit{a}_{\mathit{i}}) \middle| < \mathit{\varepsilon} \right.\]

这是 uniform continuity 的一个加强版本. uniform continuity 对应的是 \(\mathit{n} = 1\) 的情况, 即: \textbf{任意的区间上, 只要长度足够小, 就能保证 \(\mathit{F}_{\mathit{X}}\) 在这个区间上的增量足够}小; 而 absolutely continuity 更严格: \textbf{即便是任意数量的分散的区间, 只要总长度足够小, 就能保证 \(\mathit{F}_{\mathit{X}}\) 在这些 interval 上的增量的总和足够小.}

这个定义有点麻烦, 不过存在等价的条件. 为此我们要 recall 一系列定义与结论, 详细请见: \href{https://qiulinfan.github.io/mathnotes-measure_theory/12-differentiation_on_real_spaces.html}{notes on measure theory, module 12: differentiation on \(\mathbb{R}\)}:

\begin{itemize}
\item
  给定一个 function \(\mathit{F}:\mathbb{R}\rightarrow\mathbb{C}\), 我们定义它的 \textbf{total variation function} \(\mathit{V}_{\mathit{F}}:\mathbb{R}\rightarrow\lbrack 0, + \infty\rbrack\) 为:

  \[\mathit{T}_{\mathit{F}}(\mathit{x}) := \sup\left\{ \sum\limits_{\mathit{i} = 1}^{\mathit{n}} \middle| \mathit{F}(\mathit{x}_{\mathit{i}}) - \mathit{F}(\mathit{x}_{\mathit{i} - 1}) \mid - \infty < \mathit{x}_{0} < \cdots < \mathit{x}_{\mathit{n}} \right\}\]

  ~
\item
  如果 \(\mathit{T}_{\mathit{F}}(\infty) < \infty\) (注意这是单调递增函数, 因而即对于任意的 \(\mathit{x}\) 都有 \(\mathit{T}_{\mathit{F}}(\mathit{x}) < \infty\)), 那么我们称 \(\mathit{F}\) 是一个 \textbf{function of bounded variation}, 写作 \(\mathit{F} \in \mathit{B}\mathit{V}\).
\item
  如果 \(\mathit{F} \in \mathit{B}\mathit{V}\), 且 \(\mathit{F}\) 是 right-continuous 的, 且 \(\mathit{F}( - \infty) = 0\), 那么我们称 \(\mathit{F} \in \mathit{N}\mathit{B}\mathit{V}\). (即 normalized function of bounded variation).
\item
  \textbf{任意的 \(\mathit{F} \in \mathit{N}\mathit{B}\mathit{V}\) 都 \(\mathit{m}\)-a.e. differentiable, 且导数 \(\mathit{F}' \in \mathit{L}^{1}(\mathit{m})\)}, 并且以下三个条件等价:

  \[\begin{matrix}
   & {\quad\quad\mathit{F} \in \mathit{A}\mathit{C}} \\
   & {\Leftrightarrow\mathit{\mu}_{\mathit{F}} \ll \mathit{m}} \\
   & {\Leftrightarrow\mathit{F}(\mathit{x}) = \int_{- \infty}^{\mathit{x}}\mathit{F}'(\mathit{t})\mathit{d}\mathit{t}\quad\forall\mathit{x} \in \mathbb{R}}
  \end{matrix}\]
\end{itemize}

\end{qlremark}

这几行字浓缩了 measure theory 的 differentiation theory 的两个星期的内容\ldots{} 具体见上文 notes 链接, 都有详细证明. 而我们这里就利用起最后这一条结论. 首先, 我们 state 一件事:

\begin{lemma}{\kn{cdf 一定是 $\mathit{N}\mathit{B}\mathit{V}$ 的}}

任意的 random variable 的 cdf \(\mathit{F}_{\mathit{X}}\) 都是一个 \(\mathit{N}\mathit{B}\mathit{V}\) function.

\end{lemma}

\begin{proof}

依旧见 notes, 有一个关键 lemma: \(\mathit{F} \in \mathit{B}\mathit{V}\) 当且仅当它可以表示为两个 monotone increasing function 的差. 而 random variable 的 cdf \(\mathit{F}_{\mathit{X}}\) 自身是一个 non-decreasing function, 因而首先 \(\mathit{F}_{\mathit{X}} \in \mathit{B}\mathit{V}\).\\
其次, 由 \ref{thm-02-random-variables-distribution-function} 我们知道, \(\mathit{F}_{\mathit{X}}\) 是 right-continuous 的, 且 \(\mathit{F}_{\mathit{X}}( - \infty) = 0\). 因而 \(\mathit{F}_{\mathit{X}} \in \mathit{N}\mathit{B}\mathit{V}\).

\end{proof}

因而我们自然得出:

\begin{theorem}{\kn{一个 random variable 是 continuous random variable 的充要条件}}

\begin{itemize}
\item
  一个 random variable \(\mathit{X}:\Omega\rightarrow\mathbb{R}\) 的 cdf 一定是 \(\mathit{m}\)-a.e. differentiable 的.
\item
  \(\mathit{X}\) 是一个 continuous random variable 当且 仅当 \(\mathbb{P}^{\mathit{X}} \ll \mathit{m}\), 即当且仅当存在一个函数 \(\mathit{f}_{\mathit{X}}:\mathbb{R}\rightarrow\lbrack 0,\infty)\), 使得 使得对于任意 Borel set \(\mathit{A} \subseteq \mathbb{R}\), 有

  \[\mathbb{P}^{\mathit{X}}(\mathit{A}) = \int_{\mathit{A}}\mathit{f}_{\mathit{X}}\mathit{d}\mathit{m}\]

  这个函数 \(\mathit{f}_{\mathit{X}}\) 即是:

  \begin{itemize}
  \item
    \textbf{distribution 对于 Lebesgue measure 的 Radon-Nikodym derivative}: \(\mathit{f}_{\mathit{X}} = \frac{\mathit{d}\mathbb{P}^{\mathit{X}}}{\mathit{d}\mathit{m}}\).
  \item
    \textbf{cdf \(\mathit{F}_{\mathit{X}}\) 的 \(\mathit{m}\)-a.e.导数 \(\mathit{f}_{\mathit{X}} = \mathit{F}_{\mathit{X}}'\).}
  \end{itemize}

  我们也称这个 \(\mathit{f}_{\mathit{X}}\) 是 continuous random variable \(\mathit{X}\) 的 \textbf{probability density function (pdf).}
\end{itemize}

\end{theorem}

\begin{proof}

即 remark 的最后一条结论. 由于 \(\mathit{F}_{\mathit{X}}\) 是一个 \(\mathit{N}\mathit{B}\mathit{V}\) function, 因而 \(\mathit{F}_{\mathit{X}}\) 是 \(\mathit{m}\)-a.e. differentiable 的.\\
并且, \(\mathit{F}_{\mathit{X}} \in \mathit{A}\mathit{C}\) iff 它满足 FTC of Lebesgue integral.

\end{proof}

\begin{qlremark}{Remark}

因而, 在简化的教材中, 通常会直接定义: 一个 random variable \(\mathit{X}\) 是 continuous random variable, 如果存在一个 \(\mathit{f}_{\mathit{X}}\) 使得 \(\mathit{F}_{\mathit{X}}(\mathit{x}) = \int_{- \infty}^{\mathit{x}}\mathit{f}(\mathit{t})\mathit{d}\mathit{t}\) for all \(\mathit{x}\). 这里其实是跳过了我们刚才的所有步骤, 直接使用了这个等价条件. 因为这个条件也等价于:

\[\mathit{F}_{\mathit{X}}(\mathit{x}) = \int_{- \infty}^{\mathit{x}}\mathit{F}_{\mathit{X}}'(\mathit{t})\mathit{d}\mathit{t}\quad\forall\mathit{x} \in \mathbb{R}\]

\end{qlremark}

\begin{definition}{\kn{probability density function (pdf)}}

对于 continuous random variable \(\mathit{X}\), 其 probability density function 定义为 \(\mathit{F}_{\mathit{X}}\) 的 (\(\mathit{m}\)-a.e.)导数, 即

\[\mathit{f}_{\mathit{X}}(\mathit{x}) := \mathit{F}_{\mathit{X}}'(\mathit{x}),\quad\mathit{m}\ \text{-a.e.}\ \mathit{x} \in \mathbb{R}\]

\end{definition}

我们最后梳理一下.

\begin{itemize}
\item
  \textbf{任意的 random variable} 的 cdf \(\mathit{F}_{\mathit{X}}\) 都是一个 \(\mathit{N}\mathit{B}\mathit{V}\) function, 因而\textbf{一定 a.e. 可导}. 但是它的导数 \(\mathit{F}_{\mathit{X}}'\) 的积分不一定返回原函数.
\item
  \textbf{如果一个 \(\mathit{m}\)-a.e. 导数 \(\mathit{F}_{\mathit{X}}'\) 满足 \(\mathit{F}_{\mathit{X}}(\mathit{x}) = \int_{- \infty}^{\mathit{x}}\mathit{F}_{\mathit{X}}'(\mathit{t})\mathit{d}\mathit{t}\) 对所有 \(\mathit{x} \in \mathbb{R}\) 成立, 即导数的积分返回原函数, 那么 \(\mathit{X}\) 就是一个 continuous random variable}, 而我们定义 \(\mathit{f}_{\mathit{X}} := \mathit{F}_{\mathit{X}}'\) 为 \(\mathit{X}\) 的 probability density function.
\end{itemize}

以及, 显然 pdf 有以下性质:

\begin{proposition}{\kn{probability density function 的性质}}

对于 continuous random variable \(\mathit{X}\) with pdf \(\mathit{f}_{\mathit{X}}\),

\begin{itemize}
\item
  \(\mathit{f}_{\mathit{X}}(\mathit{x}) \geq 0\) for \(\mathit{m}\)-a.e. \(\mathit{x} \in \mathbb{R}\). (因为 \(\mathit{F}_{\mathit{X}}\) 是 non-decreasing 的.)
\item
  \(\int_{- \infty}^{+ \infty}\mathit{f}_{\mathit{X}}(\mathit{x})\mathit{d}\mathit{x} = 1\). (因为 \(\mathit{F}_{\mathit{X}}(\infty) = 1\).)
\item
  对于任意 Borel set \(\mathit{A}\), 有

  \[\mathbb{P}(\mathit{X} \in \mathit{A}) = \mathbb{P}^{\mathit{X}}(\mathit{A}) = \int_{\mathit{A}}\mathit{f}_{\mathit{X}}(\mathit{x})\mathit{d}\mathit{x}\]
\item
  对于 \(\mathit{m}\)-a.e. \(\mathit{x} \in \mathbb{R}\),

  \[\mathit{f}_{\mathit{X}}(\mathit{x}) = \lim\limits_{\mathit{\varepsilon}\rightarrow 0^{+}}\frac{\mathbb{P}^{\mathit{X}}(\lbrack\mathit{x},\mathit{x} + \mathit{\varepsilon}))}{\mathit{\varepsilon}} = \lim\limits_{\mathit{\varepsilon}\rightarrow 0^{+}}\frac{\mathit{F}_{\mathit{X}}(\mathit{x} + \mathit{\varepsilon}) - \mathit{F}_{\mathit{X}}(\mathit{x})}{\mathit{\varepsilon}}\]
\end{itemize}

\end{proposition}

\subsection{singular continuous random variable 与 Lebesgue Decomposition Theorem}\label{singular-continuous-random-variable-ux4e0e-lebesgue-decomposition-theorem}

对于 continuous random variable 的定义, 有些地方会简化定义为: continuous random variable 就是指其 cdf 是一个 continuous function.

但是我们现在知道, 这个定义是错的. 首先, 标准定义中的 absolutely continuous 要强于 continuous (它能推导出 a.e. differentiable); 而且, \textbf{这个更强的定义是 necessary 的}.

因为当提及 continuous random variable 时, 通常意思是它存在一个 pdf. 而存在一个 pdf 即意味着它满足 FTC of Lebesgue integral, 也就等价于 \(\mathit{F}_{\mathit{X}}\) 一定是一个 absolutely continuous function. 而仅仅 continuous 什么都无法保证.

我们这里给出一个 counterexample: Cantor distribution. 它的 cdf 是 continuous 的, 但是它并没有能够满足 FTC of Lebesgue integral (返回原函数) 的 a.e. derivative, 因而没有 pdf.

\begin{example}[(\textbf{Cantor distribution})]

令 \(\mathit{X}_{\mathit{n}}\) be 独立同分布 (i.i.d.) 的随机变量 \(\mathit{P}(\mathit{X}_{\mathit{n}} = 0) = \mathit{P}(\mathit{X}_{\mathit{n}} = 2) = 1/2\). 然后定义:

\[\mathit{X} := \sum\limits_{\mathit{n} = 1}^{\infty}\frac{\mathit{X}_{\mathit{n}}}{3^{\mathit{n}}}\]

它刻画的是:

\begin{itemize}
\item
  把 \(\lbrack 0,1\rbrack\) 三等分, 去掉中间的部分, 然后在剩下的两部分中, 以 \(1/2\) 的概率选择左边的区间, 以 \(1/2\) 的概率选择右边的区间;
\item
  在选定的区间中, 再把它三等分, infinitely 重复这一过程.
\item
  最后, 这一行为会收敛于 Cantor set 上的一个点.
\end{itemize}

具体而言:

\begin{itemize}
\item
  第 1 步 (\(\mathit{n} = 1\)) 我们查看 \(\frac{\mathit{X}_{1}}{3^{1}}\), 如果 \(\mathit{X}_{1} = 0\), 你选择了左边的区间 \(\lbrack 0,1/3\rbrack\); 如果 \(\mathit{X}_{1} = 2\), 你选择了右边的区间 \(\lbrack 2/3,1\rbrack\).
\item
  第 2 步 (\(\mathit{n} = 2\)): 在第 1 步选定的区间内，查看 \(\frac{\mathit{X}_{2}}{3^{2}} = \frac{\mathit{X}_{2}}{9}\) 如果 \(\mathit{X}_{2} = 0\), 你在当前小区间内选择了左侧的 1/3; 如果 \(\mathit{X}_{2} = 2\), 你在当前小区间内选择了右侧的 1/3.
\item
  \(\cdots\)
\end{itemize}

从而, \(\mathit{X}\) 的 range 是 Cantor set, recall: 这是一个 measure zero 的 uncountable set, 其 cardinality 是 continuum (与 \(\mathbb{R}\) 等势).

\[\mathit{C} = \left\{ {\mathit{x} \in \lbrack 0,1\rbrack \mid \mathit{x} = (0.\mathit{a}_{1}\mathit{a}_{2}\mathit{a}_{3}\ldots)_{3},\text{where}\ \mathit{a}_{\mathit{i}} \in \left\{ {0,2} \right\}} \right\}\]

它也表示了: 三进制小数中, 所有完全不包含数字 1 的小数的集合.

由于其 uncountability, 对于任意单点 \(\mathit{x}\) 都有 \(\mathbb{P}(\mathit{X} = \mathit{x}) = 0\), \textbf{因而 \(\mathit{F}_{\mathit{X}}\) 在 \(\mathbb{R}\) 上处处连续}; 并且容易证明: \textbf{\(\mathit{F}_{\mathit{X}}\) 在 \(\mathit{x}\) 处的导数 \(\mathit{F}_{\mathit{X}}'(\mathit{x}) = 0\)}. (因为对于 \(\mathit{C}^{\mathit{c}}\) 中的任意被挖掉的区间, \(\mathit{F}_{\mathit{X}}\) 在这个区间上是 constant 的).

然而, \(\mathit{F}_{\mathit{X}}' = 0\) 却\textbf{不满足 FTC of Lebesgue integral}. 比如反例: \(\mathit{F}_{\mathit{X}}(1) = 1\), 但是

\[\int_{- \infty}^{1}\mathit{F}_{\mathit{X}}'(\mathit{t})\mathit{d}\mathit{t} = 0\]

因而 \(\mathit{X}\) 没有 pdf, 也就不是 continuous random variable.

\end{example}

对于 Cantor distribution 这个例子, 我们称这样的 random variable 是一个 \textbf{singular continuous random variable}:

\begin{definition}{\kn{singular continuous random variable}}

对于 random variable \(\mathit{X}:\Omega\rightarrow\mathbb{R}\), 如果 \(\mathit{X}\) 的 cdf \(\mathit{F}_{\mathit{X}}\) 是一个 continuous function, 并且 \(\mathbb{P}^{\mathit{X}}\bot\mathit{m}\) (即存在一个 measure zero 的 Borel set \(\mathit{N}\) 使得 \(\mathbb{P}^{\mathit{X}}(\mathit{N}) = 1\), 也等价于 \(\mathit{F}_{\mathit{X}}\) 的导数 \(\mathit{F}_{\mathit{X}}' = 0\) a.e.), 那么我们称 \(\mathit{X}\) 是一个 singular continuous random variable.

\end{definition}

\begin{qlremark}{Remark}

singular continuous random variable 的意思是, 在 \(\mathbb{R}\) 上, \(\mathit{X}\) 的值 a.s. 落在一个 measure zero 的 set 上, 也就是说它的 cdf 大部分是平的, 只有在 measure zero 的 set 上才有 jump.

\end{qlremark}

这个例子向我们说明了: 除了 discrete random variable 和 continuous random variable 以外, 还有其他的奇异的 random variables. 而下面这一定理刻画了任意的 random variable 的结构:

\begin{theorem}{\kn{Lebesgue Decomposition Theorem}}

对于任意 random variable \(\mathit{X}:\Omega\rightarrow\mathbb{R}\), 其 distribution \(\mathbb{P}^{\mathit{X}}\) 可以唯一分解为三个 mutually singular 的 measure 的和:

\[\mathbb{P}^{\mathit{X}} = \mathbb{P}_{\mathit{a}\mathit{c}}^{\mathit{X}} + \mathbb{P}_{\mathit{s}\mathit{c}}^{\mathit{X}} + \mathbb{P}_{\mathit{d}}^{\mathit{X}}\]

其中 \(\mathbb{P}_{\mathit{a}\mathit{c}}^{\mathit{X}} \ll \mathit{m}\) 是 absolutely continuous measure, \(\mathbb{P}_{\mathit{s}\mathit{c}}^{\mathit{X}} \perp \mathit{m}\) 是一个 singular continuous measure, \(\mathbb{P}_{\mathit{d}}^{\mathit{X}} \perp \mathit{m}\) 是 discrete measure.

\end{theorem}

\begin{proof}

TODO.

\end{proof}

\begin{qlremark}{Remark}

我们也可以这样展开:

\[\mathit{F}_{\mathit{X}}(\mathit{x}) = \int_{- \infty}^{\mathit{x}}\mathit{f}_{\mathit{a}\mathit{c}}(\mathit{t})\,\mathit{d}\mathit{t} + \sum\limits_{\mathit{x}_{\mathit{i}} \leq \mathit{x}}\mathit{p}_{\mathit{d}}(\mathit{x}_{\mathit{i}}) + \mathit{F}_{\mathit{s}\mathit{c}}(\mathit{x})\]

其中 \(\mathit{f}_{\mathit{a}\mathit{c}}\) 是 \(\mathit{X}_{\mathit{a}\mathit{c}}\) 的 pdf, \(\mathit{p}_{\mathit{d}}\) 是 \(\mathit{X}_{\mathit{d}}\) 的 pmf, 而 \(\mathit{F}_{\mathit{s}\mathit{c}}\) 由于其奇异性是无法分解的.

\end{qlremark}

\section{expectation and variance}\label{expectation-and-variance}

Distribution 是一个 random variable 的完整刻画, 而这一节我们来谈一谈 expectation 和 variance, 这是一个 random variable 的两个重要的数值特征: 它的均值和离散程度.

最后我们会谈一谈\textbf{\(\mathit{L}^{2}(\Omega,\mathcal{F},\mathbb{P})\) 空间}: all square-integrable random variables; 以及它的一个重要 subspace \(\mathit{L}_{0}^{2}(\Omega,\mathcal{F},\mathbb{P})\): 所有 square-integrable 的 zero-centered random variables. 它们都是由 random variables 组成的 Hilbert space, 并且在其中, expectation, variance, covariance 都会得到一个非常自然的 interpretation, 展现这个空间的几何结构.

\subsection{expectation and variance 的 definition 和计算方法}\label{expectation-and-variance-ux7684-definition-ux548cux8ba1ux7b97ux65b9ux6cd5}

随机变量的 \textbf{expectation} 是它 w.r.t. 它所在概率空间 prob measure \(\mathbb{P}\) 的积分, 表示它的值的 prob-weighted average;

而其 \textbf{variance} 是 \((\mathit{X} - \mathit{E}(\mathit{X}))^{2}\) 这个 induced RV w.r.t. \(\mathbb{P}\) 的积分, 表示 \(\mathit{X}\) 离它的值的 weighted average \(\mathbb{E}(\mathit{X})\)的聚集程度:

\begin{definition}{\kn{expectation and variance of random variable}}

对于 random variable \(\mathit{X}:\Omega\rightarrow\mathbb{R}\), 其 expectation 定义为

\[\mathbb{E}\lbrack\mathit{X}\rbrack := \int_{\Omega}\mathit{X}\mathit{d}\mathbb{P}\]

其 variance 定义为

\[\text{Var}(\mathit{X}) := \mathbb{E}\lbrack(\mathit{X} - \mathbb{E}\lbrack\mathit{X}\rbrack)^{2}\rbrack = \int_{\Omega}(\mathit{X} - \mathbb{E}\lbrack\mathit{X}\rbrack)^{2}\mathit{d}\mathbb{P}\]

\end{definition}

\begin{qlremark}{Remark}

\(\int_{\Omega}\mathit{d}\mathbb{P} = \mathbb{P}(\Omega) = 1\), 因而 \(\mathbb{E}\lbrack\mathit{X}\rbrack\) 可以看作是 \(\mathit{X}\) 的值的 weighted average, 其中权重就是每个 sample point 在这个原 prob space 上的概率.

Note that: \(\mathbb{E}\lbrack\mathit{X}\rbrack\) 未必是 finite 的. \(\mathbb{E}\lbrack\mathit{X}\rbrack < \infty\Leftrightarrow\mathit{X} \in \mathit{L}^{1}(\mathbb{P})\). \(\text{Var}(\mathit{X}) < \infty\Leftrightarrow\mathit{X} \in \mathit{L}^{2}(\mathbb{P})\)

\end{qlremark}

要计算 discrete random vairable 的 expectation 和 variance, 直接 by def, sum 就可以:

\begin{proposition}{\kn{discrete random variable 的 expectation 和 variance}}

对于 discrete random variable \(\mathit{X}\) with pmf \(\mathit{p}_{\mathit{X}}\),

\[\mathbb{E}\lbrack\mathit{X}\rbrack = \sum\limits_{\mathit{x}}\mathit{x}\mathit{p}_{\mathit{X}}(\mathit{x}),\quad\text{Var}(\mathit{X}) = \sum\limits_{\mathit{x}}(\mathit{x} - \mathbb{E}\lbrack\mathit{X}\rbrack)^{2}\mathit{p}_{\mathit{X}}(\mathit{x})\]

\end{proposition}

非常简单.

但是对于 continuous random variable 而言, 怎么求呢? 这里有一个 w.r.t. prob measure \(\mathbb{P}\) 的积分, 有点难计算. 我们更希望计算 w.r.t. Lebesgue measure \(\mathit{m}\) 的积分, 这样就可以用一些经典的方法来算它了.

为了实现这个目标, 我们首先需要一个工具来过渡一下:

\begin{lemma}{\kn{用 distribution ($\mathit{X}$ 的 push-forward measure) 计算
expectation}}

对于 random variable \(\mathit{X}:\Omega\rightarrow\mathbb{R}\), 如果 \(\mathit{X}\) integrable, 那么其 expectation

\[\mathbb{E}\lbrack\mathit{X}\rbrack = \int_{\mathbb{R}}\mathit{x}\mathit{d}\mathbb{P}^{\mathit{X}}(\mathit{x})\]

\end{lemma}

\begin{proof}

这个结论是 Lebesgue integral 的 change of variable formula 的一个 special case. 注意, \(\mathit{X} = \text{id} \circ \mathit{X}\) (\(\text{id}:\mathbb{R}\rightarrow\mathbb{R}\) 是恒等函数).

By change of variable formula,

\[\int_{\Omega}\mathit{X}\mathit{d}\mathbb{P} = \int_{\Omega}(\text{id} \circ \mathit{X})\mathit{d}\mathbb{P} = \int_{\mathbb{R}}\text{id}(\mathit{x})\mathit{d}\mathbb{P}^{\mathit{X}}(\mathit{x}) = \int_{\mathbb{R}}\mathit{x}\mathit{d}\mathbb{P}^{\mathit{X}}(\mathit{x})\]

(如果不用 change of variable theorem, 也可以证明. 考虑 simple function 的 case, 显然; 然后对 simple function 自然成立, 然后用 monotone convergence theorem 推广到一般情况.)

//TODO: 从 diffeomorphism version 的 change of variable formula 证明 push-forward version 的 change of variable formula, 然后 apply 这个结论. 参考: \href{https://en.wikipedia.org/wiki/Change_of_variables}{change of variable}

\end{proof}

\begin{proposition}{\kn{continuous random variable 的 expectation 和 variance}}

对于 continuous random variable \(\mathit{X}\) with pdf \(\mathit{f}_{\mathit{X}}\), 如果 \(\mathbb{E}\lbrack\mathit{X}\rbrack\) 和 \(\text{Var}(\mathit{X})\) 都是 finite 的, 那么

\[\mathbb{E}\lbrack\mathit{X}\rbrack = \int_{- \infty}^{+ \infty}\mathit{x}\mathit{f}_{\mathit{X}}(\mathit{x})\mathit{d}\mathit{x},\quad\text{Var}(\mathit{X}) = \int_{- \infty}^{+ \infty}(\mathit{x} - \mathbb{E}\lbrack\mathit{X}\rbrack)^{2}\mathit{f}_{\mathit{X}}(\mathit{x})\mathit{d}\mathit{x}\]

并且, 对于任意的 measurable function \(\mathit{g}:\mathbb{R}\rightarrow\mathbb{R}\), 都有

\[\mathbb{E}\lbrack\mathit{g}(\mathit{X})\rbrack = \int_{- \infty}^{+ \infty}\mathit{g}(\mathit{x})\mathit{f}_{\mathit{X}}(\mathit{x})\mathit{d}\mathit{x}\]

\end{proposition}

\begin{proof}

由于 \(\mathbb{E}\lbrack\mathit{X}\rbrack = \int_{\mathbb{R}}\mathit{x}\mathit{d}\mathbb{P}^{\mathit{X}}(\mathit{x})\), 而对于 ctn random variable,

\[\mathit{d}\mathbb{P}^{\mathit{X}}(\mathit{x}) = \mathit{f}_{\mathit{X}}(\mathit{x})\mathit{d}\mathit{\lambda}(\mathit{x})\]

从而得证.

\end{proof}

\begin{proposition}{\kn{expectation 的性质}}

令 \(\mathit{X},\mathit{Y}:\Omega\rightarrow\mathbb{R}\) 为 \textbf{integrable} random variables, 则

\begin{itemize}
\item
  \textbf{linearity of expectation}: \(\mathbb{E}\lbrack\mathit{a}\mathit{X} + \mathit{b}\mathit{Y}\rbrack = \mathit{a}\mathbb{E}\lbrack\mathit{X}\rbrack + \mathit{b}\mathbb{E}\lbrack\mathit{Y}\rbrack\) for any \(\mathit{a},\mathit{b} \in \mathbb{R}\).
\item
  \textbf{monotonicity of expectation}: 如果 \(\mathit{X} \leq \mathit{Y}\) a.s., 那么 \(\mathbb{E}\lbrack\mathit{X}\rbrack \leq \mathbb{E}\lbrack\mathit{Y}\rbrack\).
\item
  \textbf{absolute value}: \(\left. \mathbb{E}\lbrack\mathit{X}\rbrack \leq \mathbb{E}\lbrack \middle| \mathit{X} \middle| \rbrack \right.\)
\end{itemize}

\end{proposition}

这三条性质都容易验证, 直接 follow from linearity of Lebesgue integral, monotonicity of Lebesgue integral, triangle inequality of Lebesgue integral.

\begin{proposition}{\kn{computing variance}}

对于 random variable \(\mathit{X}\) with finite expectation,

\[\text{Var}(\mathit{X}) = \mathbb{E}\lbrack\mathit{X}^{2}\rbrack - (\mathbb{E}\lbrack\mathit{X}\rbrack)^{2}\]

\end{proposition}

By def, 容易验证.

\subsection{\texorpdfstring{\(\mathit{L}_{0}^{2}(\Omega,\mathcal{F},\mathbb{P})\) 与 covariance and correlation as inner product and cosine}{\textbackslash mathit\{L\}\_\{0\}\^{}\{2\}(\textbackslash Omega,\textbackslash mathcal\{F\},\textbackslash mathbb\{P\}) 与 covariance and correlation as inner product and cosine}}\label{mathitl_02omegamathcalfmathbbp-ux4e0e-covariance-and-correlation-as-inner-product-and-cosine}

我们 recall measure theory 中的一丢丢 functional analysis 的内容:

\[\mathit{L}^{2}(\Omega,\mathcal{F},\mathbb{P}) = \left\{ {\mathit{X}:\Omega\rightarrow\mathbb{R} \mid \mathit{X}\ \text{is measurable(即 RV)) and}\ \mathbb{E}\lbrack\mathit{X}^{2}\rbrack < \infty} \right\}/ \sim\]

where \(\sim\) 表示 almost sure equality 的等价类, 即

\[\mathit{X} \sim \mathit{Y}\Leftrightarrow\mathbb{P}(\left\{ {\mathit{\omega} \in \Omega:\mathit{X}(\mathit{\omega}) = \mathit{Y}(\mathit{\omega})} \right\}) = 1\]

即把所有 a.s. 相等的 random variables 都看作相等.

我们知道, 这是一个 \textbf{Hilbert space}, 即一个 complete inner product space, 其中

\begin{itemize}
\item
  两个 random variables as vectors 的 inner product 定义为: \(\langle\mathit{X},\mathit{Y}\rangle := \mathbb{E}\lbrack\mathit{X}\mathit{Y}\rbrack\) 即它们的 second
\item
  一个 random variable 的 norm 即: \(\parallel \mathit{X} \parallel := \sqrt{\langle\mathit{X},\mathit{X}\rangle} = \sqrt{\mathbb{E}\lbrack\mathit{X}^{2}\rbrack}\) 即它的 second moment 的 square root. 它表示了这个 random variable 的整体离散程度.
\end{itemize}

虽然这个空间已经有一定的 information geometry 了, 但是我们实际上希望能够 用一个结构来描述两个 random variables 之间的线性相关性:

\begin{itemize}
\tightlist
\item
  任意取一个 \(\mathit{\omega}\), \(\mathit{X}(\mathit{\omega})\) 较大 (距离它的 expectation 较远) 的时候, \(\mathit{Y}(\mathit{\omega})\) 是否也通常较大 (距离它的 expectation 较远)? \(\mathit{X}(\mathit{\omega})\) 较小的时候, \(\mathit{Y}(\mathit{\omega})\) 是否也通常较小?
\end{itemize}

而 \(\mathit{L}^{2}(\Omega,\mathcal{F},\mathbb{P})\) 这个空间作为一个 Hilbert space 时, 两个 random variables 之间的距离其实是它们作为函数的相似性而不是它们变化趋势 (波动) 的相似性.

比如考虑 \(\mathit{X} = \mathit{x}\) 和 \(\mathit{Y} = \mathit{x} + 5\), 它们的线性相关性其实是 1, 即变化趋势完全相同, 但是 它们在 \(\mathit{L}^{2}(\lbrack 0,4\rbrack,\mathcal{F},\mathbb{P})\) 中的 cosine similarity 却是

\[\frac{\langle\mathit{X},\mathit{Y}\rangle}{\parallel \mathit{X} \parallel \parallel \mathit{Y} \parallel} = \frac{46/3}{\sqrt{(16/3)(151/3)}} = \frac{46}{\sqrt{2416}} \approx 0.936\]

我们希望的是这个 cosine similarity 应当是 1. 即, 应该忽略掉绝对数值, 而是考虑这两个 random variables 的变化趋势.

因而 solution: normalize 每个 random variable, 把它们移动到 均值为 0 的位置! 即: 把每个 \(\mathit{X}\) center 为 \(\mathit{X} - \mathbb{E}\lbrack\mathit{X}\rbrack\). 这就是:

\begin{definition}{\kn{zero-centered square integrable random variables
$\mathit{L}_0^2(\Omega,\mathcal{F},\mathbb{P})$}}

我们定义:

\[\mathit{L}_{0}^{2}(\Omega,\mathcal{F},\mathbb{P}) := \left\{ {\mathit{X} \in \mathit{L}^{2}(\Omega,\mathcal{F},\mathbb{P}) \mid \mathbb{E}\lbrack\mathit{X}\rbrack = 0} \right\}\]

注意, 这个空间是 \(\mathit{L}^{2}(\Omega,\mathcal{F},\mathbb{P})\) 的一个 subspace.

\end{definition}

并且, \textbf{我们可以通过把每个 \(\mathit{X}\) center 为 \(\mathit{X} - \mathbb{E}\lbrack\mathit{X}\rbrack\) 的行为, 把 \(\mathit{X}\) 从 \(\mathit{L}^{2}(\Omega,\mathcal{F},\mathbb{P})\) 投影到 \(\mathit{L}_{0}^{2}(\Omega,\mathcal{F},\mathbb{P})\) 这个 subspace 上.}\\
而投影+度量的行为, 就是计算:

\begin{definition}{\kn{covariance}}

给定两个 random variables \(\mathit{X},\mathit{Y}\), 我们定义它们的 covariance 为

\[\text{Cov}(\mathit{X},\mathit{Y}) := \mathbb{E}\lbrack(\mathit{X} - \mathbb{E}\lbrack\mathit{X}\rbrack)(\mathit{Y} - \mathbb{E}\lbrack\mathit{Y}\rbrack)\rbrack = \mathbb{E}\lbrack\mathit{X}\mathit{Y}\rbrack - \mathbb{E}\lbrack\mathit{X}\rbrack\mathbb{E}\lbrack\mathit{Y}\rbrack\]

\end{definition}

\begin{qlremark}{Remark}

Geometric intuition of covariance: 考虑取这两个 random variables 的一些 sample points 并分别作为 \((\mathit{x},\mathit{y})\) 坐标.

\begin{itemize}
\item
  如果它们的 covariance 是正的, 那么这些 sample points 倾向于拟合为一个斜率为正的 tendency line;
\item
  如果它们的 covariance 是负的, 那么这些 sample points 倾向于拟合为一个斜率为负的 tendency line;
\item
  如果它们的 covariance 是 0, 那么这些 sample points 倾向于拟合为一个水平的 tendency line.
\end{itemize}

\begin{figure}[htbp]
\centering
\csname prob-02-random-variables-diagram-01\endcsname

\caption{正、负与零协方差对应的线性趋势}

\label{fig-prob-02-random-variables-diagram-01}

\end{figure}

但是注意:

\begin{itemize}
\item
  covariance 只捕捉 linear relationship, 即趋势是否呈直线, 不过不能捕捉 nonlinear relationship. 之后在学习 independence 的时候, 我们会知道即便 covariance 为 0, \(\mathit{X}\) 和 \(\mathit{Y}\) 仍然可能是 dependent 的. 比如, 如果 \((\mathit{X},\mathit{Y})\) 的联合分布 as a graph 是个圆圈, 那么 covariance 就是 0, 但显然这两个变量之间是有关系的.

  \begin{figure}[htbp]
  \centering
  \csname prob-02-random-variables-diagram-02\endcsname

  \caption{零协方差但不独立的圆周分布}

  \label{fig-prob-02-random-variables-diagram-02}

  \end{figure}
\item
  Covariance 的数值还包括了这两个 random variables 的 scale, 因而还不能直接用来比较不同 random variables 之间的 linear relationship 的强弱程度. 真正公平的 linear relationship 还需要 normalize by norms, 把这个关系定位到 \(\lbrack - 1,1\rbrack\) 的区间上 即等下定义的 correlation.
\end{itemize}

\end{qlremark}

我们容易发现一件事情:

\begin{theorem}{\kn{covariance is a positive-semidefinite symmetric bilinear form}}

\begin{itemize}
\item
  covariance 是一个 \textbf{symmetric bilinear form} (并且 \textbf{translation-invariant}, 忽略 translation) 的 operator:

  \[\text{Cov}(\mathit{a}\mathit{X} + \mathit{b},\mathit{c}\mathit{Y} + \mathit{d}) = \mathit{a}\mathit{c}\,\text{Cov}(\mathit{X},\mathit{Y})\]\[\text{Cov}(\mathit{X},\mathit{Y}) = \text{Cov}(\mathit{Y},\mathit{X})\]\[\text{Cov}(\mathit{X} + \mathit{Y},\mathit{Z}) = \text{Cov}(\mathit{X},\mathit{Z}) + \text{Cov}(\mathit{Y},\mathit{Z})\]
\item
  covariance 满足 positive definiteness: \(\text{Cov}(\mathit{X},\mathit{X}) = \text{Var}(\mathit{X}) \geq 0\).
\end{itemize}

\end{theorem}

\begin{proof}

By def.

\end{proof}

而当我们把空间限定在 \(\mathit{L}_{0}^{2}(\Omega,\mathcal{F},\mathbb{P})\) 上时,

\begin{theorem}{\kn{$\mathit{L}_0^2(\Omega,\mathcal{F},\mathbb{P})$ 为一个 Hilbert
space, with covariance as an inner product}}

\((\mathit{L}_{0}^{2}(\Omega,\mathcal{F},\mathbb{P}),\langle \cdot , \cdot \rangle_{\text{Cov}})\) 为一个 Hilbert space, 即: covariance 在其上是一个 inner product (positive definite symmetric bilinear form).

\end{theorem}

\begin{proof}

很显然. 因为在 \(\mathit{L}^{2}\) space 上, 一个 constant random variable 的 variance 是 0, 但是它并不是一个零函数; 而在 \(\mathit{L}_{0}^{2}\) space 上, 所有的 random variables 都是 zero-centered 的, 因而 constant random variable 就是零函数了. 从而 covariance 在 \(\mathit{L}_{0}^{2}\) space 上是 positive definite 的, 从而是 inner product.\\

\end{proof}

而在 \(\mathit{L}_{0}^{2}(\Omega,\mathcal{F},\mathbb{P})\) 上, 两个 random variables \(\mathit{X},\mathit{Y}\) 之间的夹角 \(\mathit{\theta}\) 的 cosine, 就是它们 真正的 linear relationship 的强弱程度, 因而我们把它叫做 correlation:

\begin{definition}{\kn{correlation}}

\[\mathit{\rho}(\mathit{X},\mathit{Y}) = \frac{\text{Cov}(X,Y)}{\sqrt{\text{Var}(\mathit{X})}\sqrt{\text{Var}(\mathit{Y})}}\]

\end{definition}

显然,

\[\mathit{\rho}(\mathit{X},\mathit{Y}) = \cos\mathit{\theta}\]

where \(\mathit{\theta}\) 是 \(\mathit{X},\mathit{Y}\) 之间的夹角. And we have:

\[\text{Cov}(\mathit{X},\mathit{Y}) = \mathit{\sigma}_{\mathit{X}}\mathit{\sigma}_{\mathit{Y}}\mathit{\rho}(\mathit{X},\mathit{Y})\]

从而有以下显然的结论:

\begin{proposition}{\kn{properties of correlation}}

\begin{itemize}
\item
  \(\mathit{\rho}(\mathit{X},\mathit{Y}) \in \lbrack - 1,1\rbrack\)
\item
  如果 \(\mathit{\rho}(\mathit{X},\mathit{Y}) = 1\), 那么 \(\mathit{Y} = \mathit{a}\mathit{X} + \mathit{b}\) for some \(\mathit{a} > 0\) and \(\mathit{b}\); 如果 \(\mathit{\rho}(\mathit{X},\mathit{Y}) = - 1\), 那么 \(\mathit{Y} = \mathit{a}\mathit{X} + \mathit{b}\) for some \(\mathit{a} < 0\) and \(\mathit{b}\).~
\item
  By Cauchy-Schwarz inequality, \(\left. |\text{Cov}(\mathit{X},\mathit{Y}) \middle| \leq \mathit{\sigma}_{\mathit{X}}\mathit{\sigma}_{\mathit{Y}} \right.\).
\end{itemize}

\end{proposition}

做一个总结:

{\def\LTcaptype{none} % do not increment counter
\small\setlength{\tabcolsep}{4pt}\renewcommand{\arraystretch}{1.3}
\begin{longtable}[]{@{}p{0.23\linewidth}p{0.31\linewidth}p{0.39\linewidth}@{}}
\toprule\noalign{}
& \textbf{\(\mathit{L}^{2}\) space} & \textbf{\(\mathit{L}_{0}^{2}\)} \\
\midrule\noalign{}
\endhead
\bottomrule\noalign{}
\endlastfoot
inner product \(\langle\mathit{X},\mathit{Y}\rangle\) & \(\mathbb{E}\lbrack\mathit{X}\mathit{Y}\rbrack\) & \(\text{Cov}(\mathit{X},\mathit{Y})\) (covariance) \\
norm \(\parallel \mathit{X} \parallel\) & \(\sqrt{\mathbb{E}\lbrack\mathit{X}^{2}\rbrack}\) & \(\sqrt{\text{Var}(\mathit{X})}\) (standard deviation) \\
cosine similarity & \(\frac{\mathbb{E}\lbrack\mathit{X}\mathit{Y}\rbrack}{\sqrt{\mathbb{E}\lbrack\mathit{X}^{2}\rbrack}\sqrt{\mathbb{E}\lbrack\mathit{Y}^{2}\rbrack}}\) & \(\frac{\text{Cov}(\mathit{X},\mathit{Y})}{\sqrt{\text{Var}(\mathit{X})}\sqrt{\text{Var}(\mathit{Y})}}\) (correlation) \\
\end{longtable}
}

\textbf{求两个 random variables 之间的 correlation, 就是把它们都 center 到 zero-centered 的位置 (即投影到 \(\mathit{L}_{0}^{2}\) space 上), 然后计算它们在 \(\mathit{L}_{0}^{2}\) space 上的 cosine similarity.}\\
此时我们看到 decomposition of variance:

\begin{proposition}{\kn{两个 random variables 之和的 variance}}

两个 random variables 的和的 variance 可分解成 variances 和 covariance:

\[\text{Var}(\mathit{X} + \mathit{Y}) = \text{Var}(\mathit{X}) + \text{Var}(\mathit{Y}) + 2\ \text{Cov}(\mathit{X},\mathit{Y})\]

\end{proposition}

\begin{proof}

\[\text{Var}(\mathit{X} + \mathit{Y}) = \mathbb{E}\left\lbrack {(\mathit{X} + \mathit{Y})^{2}} \right\rbrack - (\mathbb{E}\lbrack\mathit{X} + \mathit{Y}\rbrack)^{2} = \mathbb{E}\left\lbrack \mathit{X}^{2} \right\rbrack - (\mathbb{E}\lbrack\mathit{X}\rbrack)^{2} + \mathbb{E}\left\lbrack \mathit{Y}^{2} \right\rbrack - (\mathbb{E}\lbrack\mathit{Y}\rbrack)^{2} + 2(\mathbb{E}\lbrack\mathit{X}\mathit{Y}\rbrack - \mathbb{E}\lbrack\mathit{X}\rbrack\mathbb{E}\lbrack\mathit{Y}\rbrack)\]

\end{proof}

\begin{qlremark}{Remark}

它其实这就是标准的 inner product space 上的向量平方公式:

\[\parallel \mathit{X} + \mathit{Y}\underset{\mathit{L}_{0}^{2}}{\overset{2}{\parallel}} = \parallel \mathit{X}\underset{\mathit{L}_{0}^{2}}{\overset{2}{\parallel}} + \parallel \mathit{Y}\underset{\mathit{L}_{0}^{2}}{\overset{2}{\parallel}} + 2\langle\mathit{X},\mathit{Y}\rangle_{\mathit{L}_{0}^{2}}\]

\end{qlremark}

\section{discrete random variables}\label{discrete-random-variables}

Recall: discrete random variable 是指其 range a.s. 是 countable 的 random variable. (存在一个 countable set \(\mathit{S} \subseteq \mathbb{R}\) 使得 \(\mathbb{P}(\mathit{X} \in \mathit{S}) = 1\))

discrete RV 的 probability mass unction (pmf):

\[\mathit{p}_{\mathit{X}}(\mathit{x}) := \mathbb{P}(\mathit{X} = \mathit{x}),\quad\forall\mathit{x} \in \mathbb{R}\]

就是 \(\mathit{X}\) 的 distribution \(\mathit{P}^{\mathit{X}}\) 对于 counting measure of range \(\mathit{S}\) 的 Radon-Nikodym derivative:

\[\mathit{p}_{\mathit{X}} = \frac{\mathit{d}\mathbb{P}^{\mathit{X}}}{\mathit{d}\mathit{\mu}_{\text{S}}}\]

因为假设 \(\mathit{X}\) 的 range 是 \(\mathit{S} = \left\{ {\mathit{x}_{1},\mathit{x}_{2},\ldots} \right\}\), 那么对于任意 Borel set \(\mathit{B}\),

\[\mathbb{P}^{\mathit{X}}(\mathit{B}) = \sum\limits_{\mathit{x}_{\mathit{i}} \in \mathit{B}}\mathbb{P}(\mathit{X} = \mathit{x}_{\mathit{i}}) = \sum\limits_{\mathit{x}_{\mathit{i}} \in \mathit{B} \cap \mathit{S}}\mathit{p}_{\mathit{X}}(\mathit{x}_{\mathit{i}}) = \int_{\mathit{A}}\mathit{p}_{\mathit{X}}(\mathit{x})\mathit{d}\mathit{\mu}_{\text{S}}(\mathit{x})\]

下面我们给出一些经典的 discrete random variable 的例子, 以及它们的 pmf 和 cdf.

\subsection{Bernoulli distribution and Binomial distribution}\label{bernoulli-distribution-and-binomial-distribution}

\begin{example}[(\textbf{Bernoulli distribution})]

令 \(\Omega := \left\{ {\mathit{\omega}_{1},\mathit{\omega}_{2}} \right\}\), \(\mathcal{F} := 2^{\Omega}\).\\
令 \(\mathbb{P}(\left\{ \mathit{\omega}_{1} \right\}) = \mathit{p}\), \(\mathbb{P}(\left\{ \mathit{\omega}_{2} \right\}) = 1 - \mathit{p}\) 来表示这两个 singular 事件的概率.\\
令 \(\mathit{X}(\mathit{\omega}_{1}) = 1\), \(\mathit{X}(\mathit{\omega}_{2}) = 0\) 分别表示 success 和 failure.\\
那么显然可以计算:

\[\mathbb{P}(\mathit{X} = 0) = \mathbb{P}(\left\{ \mathit{\omega}_{2} \right\}) = 1 - \mathit{p},\quad\mathbb{P}(\mathit{X} = 1) = \mathbb{P}(\left\{ \mathit{\omega}_{1} \right\}) = \mathit{p}\]

我们称 \((\Omega,\mathcal{F},\mathbb{P})\) 为一个 Bernoulli probability space (它 model 了一个 Bernoulli trial, 即一个 random experiment with two possible outcomes: success 和 failure)\\
称 \(\mathit{X}:\Omega\rightarrow\left\{ {0,1} \right\}\) 这个 random variable 为一个 \textbf{Bernoulli random variable}, 并称 \(\mathit{X}\) 的 distribution \(\mathit{P}^{\mathit{X}}\) 为一个 \textbf{Bernoulli distribution}, 写作

\[\mathit{X} \sim \text{Ber}(\mathit{p})\]

这是最简单的 random variable 和 distribution 了. 它 model 的是: 比如我们 toss 一枚 biased coin, 以 \(\mathit{p}\) 的概率得到 heads (success), 以 \(1 - \mathit{p}\) 的概率得到 tails (failure).\\

\end{example}

\begin{proposition}

如果 \(\mathit{X} \sim \text{Ber}(\mathit{p})\), 那么 \(\mathbb{E}\lbrack\mathit{X}\rbrack = \mathit{p}\), \(\text{Var}(\mathit{X}) = \mathit{p}(1 - \mathit{p})\).

\end{proposition}

显然.

\begin{example}[(\textbf{Binomial distribution})]

我们 independently repeat \(\mathit{n}\) 次 Bernoulli trial, 每次 trial 的 success probability 都是 \(\mathit{p}\).\\
考虑:

\[\mathit{S} := \mathit{X}_{1} + \mathit{X}_{2} + \cdots + \mathit{X}_{\mathit{n}}\]

为 success 的总次数. 那么 \(\mathit{S}\) 是一个 discrete random variable \(\mathit{S}:\Omega^{\mathit{n}}\rightarrow\mathbb{Z}_{\geq 0}\) 容易计算出 \(\mathit{S}\) 的 pmf:

\[\mathit{p}_{\mathit{S}}(\mathit{k}) = \mathbb{P}(\mathit{S} = \mathit{k}) = \binom{\mathit{n}}{\mathit{k}}\mathit{p}^{\mathit{k}}(1 - \mathit{p})^{\mathit{n} - \mathit{k}},\quad\mathit{k} = 0,1,\ldots,\mathit{n}\]

我们称 \(\mathit{S}\) 的 distribution \(\mathit{P}^{\mathit{S}}\) 为一个 \textbf{Binomial distribution}, 写作

\[\mathit{S} \sim \text{Bin}(\mathit{n},\mathit{p})\]

它 model 的是: 比如我们 toss 一枚 biased coin \(\mathit{n}\) 次, 得到 heads 的总次数.\\

\end{example}

\begin{proposition}

如果 \(\mathit{S} \sim \text{Bin}(\mathit{n},\mathit{p})\), 那么 \(\mathbb{E}\lbrack\mathit{S}\rbrack = \mathit{n}\mathit{p}\), \(\text{Var}(\mathit{S}) = \mathit{n}\mathit{p}(1 - \mathit{p})\).

\end{proposition}

\begin{proof}

\[\mathbb{E}\lbrack\mathit{X}\rbrack = \sum\limits_{\mathit{k} = 1}^{\mathit{n}}\mathit{k} \cdot \mathbb{P}(\mathit{X} = \mathit{k}) = \sum\limits_{\mathit{k} = 1}^{\mathit{n}}\mathit{k} \cdot \binom{\mathit{n}}{\mathit{k}}\mathit{p}^{\mathit{k}}(1 - \mathit{p})^{\mathit{n} - \mathit{k}} = \mathit{p}(1 - \mathit{p})^{\mathit{n} - 1}\sum\limits_{\mathit{k} = 1}^{\mathit{n}}\mathit{k} \cdot \binom{\mathit{n}}{\mathit{k}}(\frac{\mathit{p}}{1 - \mathit{p}})^{\mathit{k} - 1}\]

由于

\[(1 + \mathit{t})^{\mathit{n}} = \sum\limits_{\mathit{k} = 0}^{\mathit{n}}\binom{\mathit{n}}{\mathit{k}}\mathit{t}^{\mathit{k}}\Longrightarrow\mathit{n}(1 + \mathit{t})^{\mathit{n} - 1} = \sum\limits_{\mathit{k} = 1}^{\mathit{n}}\mathit{k} \cdot \binom{\mathit{n}}{\mathit{k}}\mathit{t}^{\mathit{k} - 1}\]

可以得到

\[\mathbb{E}\lbrack\mathit{X}\rbrack = \mathit{p}(1 - \mathit{p})^{\mathit{n} - 1} \cdot \mathit{n}(1 + \frac{\mathit{p}}{1 - \mathit{p}})^{\mathit{n} - 1} = \mathit{n}\mathit{p}\]

然后计算得 \(\text{Var}(\mathit{X}) = \mathit{n}\mathit{p}(1 - \mathit{p})\). 但是这是比较麻烦的方法. 也可以利用 \(\mathit{X} = \sum_{\mathit{i} = 1}^{\mathit{n}}\mathit{X}_{\mathit{i}}\), 其中 \(\mathit{X}_{\mathit{i}} \sim \text{Ber}(\mathit{p})\) i.i.d 这个事实来计算, 那么直接通过 线性叠加 Bernoulli random variable 的 expectation 和 variance 就可以了.

\end{proof}

\subsection{Geometric distribution and Negative Binomial distribution}\label{geometric-distribution-and-negative-binomial-distribution}

\begin{example}[(\textbf{Geometric distribution})]

我们 perform independent Bernoulli trial, 每次 trial 的 success probability 都是 \(\mathit{p}\).\\
考虑 random variable \(\mathit{T}\) 表示第一次 success 发生的 trial number. (也就是等于: 我们一直 trial, 直到第一次 success 发生, 那么这个 trial 的 number 就是 \(\mathit{T}\)).\\
那么 \(\mathit{T}\) 是一个 discrete random variable \(\mathit{T}:\Omega^{\infty}\rightarrow\mathbb{Z}_{\geq 1}\).\\
容易计算出 \(\mathit{T}\) 的 pmf:

\[\mathit{p}_{\mathit{T}}(\mathit{k}) = \mathbb{P}(\mathit{T} = \mathit{k}) = (1 - \mathit{p})^{\mathit{k} - 1}\mathit{p},\quad\mathit{k} = 1,2,\ldots\]

我们称 \(\mathit{T}\) 的 distribution \(\mathit{P}^{\mathit{T}}\) 为一个 \textbf{Geometric distribution}, 写作

\[\mathit{T} \sim \text{Geom}(\mathit{p})\]

它 model 的是: 比如我们 toss 一枚 biased coin, 第一次得到 heads 的 trial number.\\

\end{example}

\begin{proposition}

如果 \(\mathit{X} \sim \text{Geom}(\mathit{p})\), 那么 \(\mathbb{E}\lbrack\mathit{X}\rbrack = 1/\mathit{p}\), \(\text{Var}(\mathit{X}) = (1 - \mathit{p})/\mathit{p}^{2}\).

\end{proposition}

\begin{proof}

\[\mathbb{E}\lbrack\mathit{X}\rbrack = \sum\limits_{\mathit{k} = 1}^{\infty}\mathit{k}(1 - \mathit{p})^{\mathit{k} - 1}\mathit{p} = \mathit{p}\sum\limits_{\mathit{k} = 1}^{\infty}\mathit{k}(1 - \mathit{p})^{\mathit{k} - 1} = \mathit{p}( - \frac{\mathit{d}}{\mathit{d}(1 - \mathit{p})}\sum\limits_{\mathit{k} = 0}^{\infty}(1 - \mathit{p})^{\mathit{k}}) = - \mathit{p} - \frac{\mathit{d}}{\mathit{d}\mathit{p}}\frac{1}{1 - (1 - \mathit{p})} = \frac{1}{\mathit{p}}\]

similarly, 可以计算得

\[\mathbb{E}\lbrack\mathit{X}^{2}\rbrack = \frac{2 - \mathit{p}}{\mathit{p}^{2}}\]

因而

\[\text{Var}\lbrack\mathit{X}\rbrack = \mathbb{E}\lbrack\mathit{X}^{2}\rbrack - (\mathbb{E}\lbrack\mathit{X}\rbrack)^{2} = \frac{1 - \mathit{p}}{\mathit{p}^{2}}\]

\end{proof}

\begin{example}[(\textbf{Negative Binomial distribution})]

这是 Geometric distribution 的 generalization (但不完全一样).\\
我们 perform independent Bernoulli trial, 每次 trial 的 success probability 都是 \(\mathit{p}\). (即 independent and identically distributed)\\
令 random variable \(\mathit{T}_{\mathit{r}}\) 表示第 \(\mathit{r}\) 次 success 发生之前 的 failures 的数量. (也就是等于: 我们一直 trial, 直到第 \(\mathit{r}\) 次 success 发生, 那么这个 trial 的 number 就是 \(\mathit{T}_{\mathit{r}} + \mathit{r}\), 因为 \(\mathit{T}_{\mathit{r}}\) 是 failures 的数量, 还要加上 \(\mathit{r} - 1\) 个 success. ).\\
那么 \(\mathit{T}_{\mathit{r}}\) 是一个 discrete random variable \(\mathit{T}_{\mathit{r}}:\Omega^{\infty}\rightarrow\mathbb{Z}_{\geq 0}\).\\
容易计算出 \(\mathit{T}_{\mathit{r}}\) 的 pmf:

\[\mathit{p}_{\mathit{T}_{\mathit{r}}}(\mathit{k}) = \mathbb{P}(\mathit{T}_{\mathit{r}} = \mathit{k}) = \binom{\mathit{k} + \mathit{r} - 1}{\mathit{k}}(1 - \mathit{p})^{\mathit{k}}\mathit{p}^{\mathit{r}},\quad\mathit{k} = 0,1,2,\ldots\]

这是因为:最后一次 success 的位置是固定的. 我们要在前 \(\mathit{k} + \mathit{r} - 1\) 次 trial 中选择 \(\mathit{k}\) 次作为 failures.\\
我们称 \(\mathit{T}_{\mathit{r}}\) 的 distribution \(\mathit{P}^{\mathit{T}_{\mathit{r}}}\) 为一个 \textbf{Negative Binomial distribution}, 写作

\[\mathit{T}_{\mathit{r}} \sim \text{NB}(\mathit{r},\mathit{p})\]

它 model 的是: 比如我们 toss 一枚 biased coin, 得到第 \(\mathit{r}\) 个 heads 之前会经历的 failures 的数量.\\

\end{example}

\subsection{Poisson distribution}\label{poisson-distribution}

\begin{example}[(\textbf{Poisson distribution})]

考虑这一 pmf:

\[\mathbb{P}(\mathit{X} = \mathit{k}) = \mathit{e}^{- \mathit{\lambda}}\frac{\mathit{\lambda}^{\mathit{k}}}{\mathit{k}!},\quad\mathit{k} = 0,1,2,\ldots\]

我们称 \(\mathit{X}\) 的 distribution \(\mathit{P}^{\mathit{X}}\) 为一个 \textbf{Poisson distribution}, 写作

\[\mathit{X} \sim \text{Poi}(\mathit{\lambda}),\quad\mathit{\lambda} > 0\]

以 \(\mathit{\lambda} = 3\) 为例画出 pmf 示意图:

\begin{figure}[htbp]
\centering
\csname prob-02-random-variables-diagram-03\endcsname

\caption{参数 \(\mathit{\lambda} = 3\) 的 Poisson 概率质量函数}

\label{fig-prob-02-random-variables-diagram-03}

\end{figure}

\begin{qlremark}{Remark}

当 \(\mathit{n}\rightarrow\infty\) 时, \(\text{Bin}(\mathit{n},\mathit{\lambda}/\mathit{n})\) 的 distribution 趋近于 \(\text{Poi}(\mathit{\lambda})\).\\
考虑 \(\mathit{X} \sim \mathit{B}\left( {\mathit{n},\frac{\mathit{\lambda}}{\mathit{n}}} \right)\), 对于固定的 \(\mathit{k}\):

\[\mathbb{P}(\mathit{X} = \mathit{k}) = \binom{\mathit{n}}{\mathit{k}}\left( \frac{\mathit{\lambda}}{\mathit{n}} \right)^{\mathit{k}}\left( {1 - \frac{\mathit{\lambda}}{\mathit{n}}} \right)^{\mathit{n} - \mathit{k}}\]

展开组合数 \(\binom{\mathit{n}}{\mathit{k}}\):

\[\mathbb{P}(\mathit{X} = \mathit{k}) = \frac{\mathit{n}(\mathit{n} - 1)\ldots(\mathit{n} - \mathit{k} + 1)}{\mathit{k}!} \cdot \frac{\mathit{\lambda}^{\mathit{k}}}{\mathit{n}^{\mathit{k}}} \cdot \left( {1 - \frac{\mathit{\lambda}}{\mathit{n}}} \right)^{\mathit{n}} \cdot \left( {1 - \frac{\mathit{\lambda}}{\mathit{n}}} \right)^{- \mathit{k}}\]

重组项:

\[\mathbb{P}(\mathit{X} = \mathit{k}) = \frac{\mathit{\lambda}^{\mathit{k}}}{\mathit{k}!} \cdot \underset{(1)}{\underbrace{\frac{\mathit{n}(\mathit{n}-1)\ldots(\mathit{n}-\mathit{k}+1)}{\mathit{n}^{\mathit{k}}}}} \cdot \underset{(2)}{\underbrace{\left( {1-\frac{\mathit{\lambda}}{\mathit{n}}} \right)^{\mathit{n}}}} \cdot \underset{(3)}{\underbrace{\left( {1-\frac{\mathit{\lambda}}{\mathit{n}}} \right)^{-\mathit{k}}}}\]

其中, (1) 的 limit 是 1, (2) 的 limit 是 \(\mathit{e}^{- \mathit{\lambda}}\), (3) 的 limit 是 1, 因而

\[\lim\limits_{\mathit{n}\rightarrow\infty}\mathbb{P}(\mathit{X} = \mathit{k}) = \mathit{e}^{- \mathit{\lambda}}\frac{\mathit{\lambda}^{\mathit{k}}}{\mathit{k}!}\]

这个极限行为的意义是:

\begin{itemize}
\item
  \(\mathit{\lambda}\) 表示: 在一个固定时间/空间窗口内, 一个事件出现的平均次数的观测.
\item
  我们把这段时间间分成 \(\mathit{n}\) 个小窗口, 将这个观测作为先验知识用来 model 这类事件实现的平均概率, 那么 每个小窗口内事件出现的概率是 \(\mathit{p} = \mathit{\lambda}/\mathit{n}\).
\item
  每个 \(\text{Bin}(\mathit{n},\mathit{\lambda}/\mathit{n})\) 都表示: 在时间细分成 \(\mathit{n}\) 个小窗口的情况下, 这个事件出现的次数的分布. (注意不是在时间上的分布, 而是在次数上的分布)
\item
  当 \(\mathit{n}\rightarrow\infty\) 时, 其意义: 基于连续时间的估计, 这个事件在原固定事件窗口内出现的次数的分布. 这就是 \(\text{Poi}(\mathit{\lambda})\) 的意义.
\end{itemize}

因而 Poisson distribution 的意义本质上是: \textbf{给定对一个事件在固定窗口下出现的平均次数的观测(估计), 那么这个事件在这个窗口内出现 \(\mathit{n}\) 次 (\(\mathit{n} = 0,1,2,\cdots\))的概率是多少?}\\
所以 Poisson 分布, 即是\textbf{根据一个事件在给定 windows 内发生次数的 expectation, 来估计这个事件在该 windows 内发生次数的分布的标准方法.}\\

\end{qlremark}

\begin{qlremark}{Remark}

Further remark: Poisson distribution 的用途.

\begin{itemize}
\item
  它是 model "大量级的小概率事件" 的通用模型.\\
  即一段时间内, 一件事情的发生的概率非常小 (\(\mathit{p} \ll 1\)), 但是有很多机会发生 (\(\mathit{n} \gg 1\)). 比如放射性衰变.
\item
  现实世界中, 对于一些事情, 我们想要 model 其在一段时间内发生的总次数 (\(\mathit{n}\)) 和单次成功率 (\(\mathit{p}\)) 是很困难的. 比如统计今天路过百货大楼的总人数 (\(\mathit{n}\)), 以及每个人路过百货大楼时走进去的概率 (\(\mathit{p}\)). 但是我们 很容易可以知道: \(\mathit{\lambda} = \mathit{n}\mathit{p}\) 的观测. 即: 每天平均有多少人走进百货大楼. 那么我们就可以用 Poisson distribution 来 model 这件事情在一天内发生的次数的分布.\\
\end{itemize}

\end{qlremark}

\end{example}

\begin{proposition}

如果 \(\mathit{X} \sim \text{Poi}(\mathit{\lambda})\), 那么 \(\mathbb{E}\lbrack\mathit{X}\rbrack = \mathit{\lambda}\), \(\text{Var}(\mathit{X}) = \mathit{\lambda}\).

\end{proposition}

\begin{proof}

recall \(\mathit{e}^{\mathit{x}}\) 的 taylor expansion:

\[\mathit{e}^{\mathit{x}} = \sum\limits_{\mathit{k} = 0}^{\infty}\frac{\mathit{x}^{\mathit{k}}}{\mathit{k}!}\]

因而

\[\mathbb{E}\lbrack\mathit{X}\rbrack = \sum\limits_{\mathit{k} = 0}^{\infty}\mathit{k}\mathit{e}^{- \mathit{\lambda}}\frac{\mathit{\lambda}^{\mathit{k}}}{\mathit{k}!} = \mathit{e}^{- \mathit{\lambda}}\mathit{\lambda}\sum\limits_{\mathit{k} = 1}^{\infty}\frac{\mathit{\lambda}^{\mathit{k} - 1}}{(\mathit{k} - 1)!} = \mathit{e}^{- \mathit{\lambda}}\mathit{\lambda}\mathit{e}^{\mathit{\lambda}} = \mathit{\lambda}\]

similarly, 可以计算得

\[\mathbb{E}\lbrack\mathit{X}^{2}\rbrack = \mathit{\lambda}^{2} + \mathit{\lambda}\]

因此,

\[\text{Var}(\mathit{X}) = \mathbb{E}\lbrack\mathit{X}^{2}\rbrack - (\mathbb{E}\lbrack\mathit{X}\rbrack)^{2} = \mathit{\lambda}\]

\end{proof}

\subsubsection{Poisson distribution 的 additivity}\label{poisson-distribution-ux7684-additivity}

\begin{theorem}{\kn{additivity of Poisson distribution}}

如果 \(\mathit{X} \sim \text{Poi}(\mathit{\lambda}_{1})\), \(\mathit{Y} \sim \text{Poi}(\mathit{\lambda}_{2})\) are independent, 那么

\[\mathit{X} + \mathit{Y} \sim \text{Poi}(\mathit{\lambda}_{1} + \mathit{\lambda}_{2})\]

\end{theorem}

\begin{qlremark}{Remark}

这个事实通过简单的计算即可得到.

但是值得一提的是, 我们之后会证明一个结论: 两个 independent 的 random variable 的 sum 的 distribution 是它们 convolution 的 distribution.

而 Poisson distribution 的 convolution 仍然是 Poisson distribution, 并且是 linear 的. 这就是 additivity of Poisson distribution 的本质原因.

\end{qlremark}

\section{continuous random variables}\label{continuous-random-variables}

recall: 总而言之, continuous random variable 是指其 distribution \(\mathbb{P}^{\mathit{X}}\) 对于 Lebesgue measure \(\mathit{m}\) 是 absolutely continuous 的 random variable, 等价于存在一个 pdf \(\mathit{f}_{\mathit{X}}\) 使得

\[\mathit{F}_{\mathit{X}}(\mathit{x}) = \int_{- \infty}^{\mathit{x}}\mathit{f}_{\mathit{X}}(\mathit{y})\mathit{d}\mathit{y}\]

对于任意 \(\mathit{x}\) 都成立.

而由 \(\mathbb{E}\lbrack\mathit{X}\rbrack = \int_{\Omega}\mathit{X}\mathit{d}\mathbb{P} = \int_{\mathbb{R}}\mathit{x}\mathit{d}\mathbb{P}^{\mathit{X}} = \int_{\mathbb{R}}\mathit{x}\mathit{f}_{\mathit{X}}(\mathit{x})\mathit{d}\mathit{x}\) 可得

\[\mathbb{E}\lbrack\mathit{X}\rbrack = \int_{- \infty}^{+ \infty}\mathit{x}\mathit{f}_{\mathit{X}}(\mathit{x})\mathit{d}\mathit{x},\quad\text{Var}(\mathit{X}) = \int_{- \infty}^{+ \infty}(\mathit{x} - \mathbb{E}\lbrack\mathit{X}\rbrack)^{2}\mathit{f}_{\mathit{X}}(\mathit{x})\mathit{d}\mathit{x}\]

这里还有一个额外的结论:

\begin{proposition}

对于任意的 measurable function \(\mathit{g}:\mathbb{R}\rightarrow\mathbb{R}\), 都有

\[\mathbb{E}\lbrack\mathit{g}(\mathit{X})\rbrack = \int_{- \infty}^{+ \infty}\mathit{g}(\mathit{x})\mathit{f}_{\mathit{X}}(\mathit{x})\mathit{d}\mathit{x}\]

\end{proposition}

\begin{proof}

\[\begin{matrix}
{\mathbb{E}\lbrack\mathit{g}(\mathit{X})\rbrack} & {= \int_{\Omega}\mathit{g}(\mathit{X}(\mathit{w}))\,\mathit{d}\mathbb{P}(\mathit{w})} \\
 & {= \int_{\Omega}\mathit{g} \circ \mathit{X}\,\mathit{d}\mathbb{P}} \\
 & {= \int_{\mathbb{R}}\mathit{g}\,\mathit{d}\mathbb{P}^{\mathit{X}}} \\
 & {= \int_{- \infty}^{+ \infty}\mathit{g}(\mathit{x})\mathit{f}_{\mathit{X}}(\mathit{x})\,\mathit{d}\mathit{x}}
\end{matrix}\]

\end{proof}

下面我们介绍一些常见的 continuous random variables, 以及它们的 pdf, cdf, expectation 和 variance.

\subsection{uniform distribution: 最简单的 continuous RV}\label{uniform-distribution-ux6700ux7b80ux5355ux7684-continuous-rv}

\begin{definition}{\kn{uniform distribution}}

如果 \(\mathit{f}_{\mathit{X}}(\mathit{x}) = \frac{1}{\mathit{b} - \mathit{a}}\) on \((\mathit{a},\mathit{b})\) and \(0\) elsewhere, 那么我们称 \(\mathit{X}\) 的 distribution \(\mathit{P}^{\mathit{X}}\) 为一个 \textbf{uniform distribution}, 写作

\[\mathit{X} \sim \mathit{U}(\lbrack\mathit{a},\mathit{b}\rbrack)\]

\end{definition}

对于 uniform distribution, 很容易计算得:

\[\mathit{F}_{\mathit{X}}(\mathit{x}) = \frac{\mathit{x} - \mathit{b}}{\mathit{b} - \mathit{a}},\quad\mathbb{E}\lbrack\mathit{X}\rbrack = \frac{\mathit{a} + \mathit{b}}{2},\quad\text{Var}(\mathit{X}) = \frac{(\mathit{b} - \mathit{a})^{2}}{12}\]

\begin{example}[(uniform distribution 的 pdf 和 cdf)]

考虑 \(\mathit{X} \sim \mathit{U}(\lbrack\mathit{a},\mathit{b}\rbrack)\), 其 pdf 和 cdf 的图像如下:

\begin{figure}[htbp]
\centering
\csname prob-02-random-variables-diagram-04\endcsname

\caption{区间 \(\left\lbrack {\mathit{a},\mathit{b}} \right\rbrack\)
上均匀分布的概率密度函数}

\label{fig-prob-02-random-variables-diagram-04}

\end{figure}

\begin{figure}[htbp]
\centering
\csname prob-02-random-variables-diagram-05\endcsname

\caption{区间 \(\left\lbrack {\mathit{a},\mathit{b}} \right\rbrack\)
上均匀分布的累积分布函数}

\label{fig-prob-02-random-variables-diagram-05}

\end{figure}

\end{example}

\subsection{exponential distribution: geometric distribution 的 continuous version}\label{exponential-distribution-geometric-distribution-ux7684-continuous-version}

Exponential random variable model 的是一盏灯的 remaining lifetime. 它 assume: 任意时间, 这盏灯有一个 constant 的磨损速率.

假设它的磨损速率是 \(\mathit{\lambda} = \mathit{\lambda}_{\mathit{X}}(\mathit{t})\) for all \(\mathit{t}\), 那么在任意时间 \(\mathit{t}\) 上:

\[\mathit{\lambda}_{\mathit{X}}(\mathit{t}) = \lim\limits_{h\rightarrow 0^{+}}\frac{\mathit{P}(\mathit{X} \in (\mathit{t},\mathit{t} + h\rbrack)}{h} = \lim\limits_{h\rightarrow 0^{+}}\frac{\mathbb{P}(\mathit{t} \leq \mathit{X} \leq \mathit{t} + h)}{h} \cdot \frac{1}{\mathbb{P}(\mathit{X} > \mathit{t})} = \frac{\mathit{f}_{\mathit{X}}(\mathit{t})}{1 - \mathit{F}_{\mathit{X}}(\mathit{t})}\]

即 ODE:

\[\mathit{\lambda} = \frac{\mathit{F}_{\mathit{X}}'(\mathit{t})}{1 - \mathit{F}_{\mathit{X}}(\mathit{t})}\]

with initial condition \(\mathit{F}_{\mathit{X}}(0) = 0\).

我们知道这个 ODE 的 solution 是:

\[\begin{matrix}
{\mathit{F}_{\mathit{X}}(\mathit{t}) = \{ 1 - \mathit{e}^{- \mathit{\lambda}\mathit{t}},\quad\mathit{t} > 0} \\
{0,\quad\mathit{t} \leq 0,\quad\mathit{f}_{\mathit{X}}(\mathit{t}) = \{\mathit{\lambda}\mathit{e}^{- \mathit{\lambda}\mathit{t}},\quad\mathit{t} > 0} \\
{0,\quad\mathit{t} \leq 0}
\end{matrix}\]

\begin{definition}{\kn{exponential distribution}}

我们称 \(\mathit{X}\) 为一个 \textbf{exponential random variable} with parameter \(\mathit{\lambda}\), 如果它的 pdf is given by:

\[\begin{matrix}
{\mathit{f}(\mathit{x}) = \{\mathit{\lambda}\mathit{e}^{- \mathit{\lambda}\mathit{x}},\quad\mathit{x} > 0} \\
{0,\quad\mathit{x} \leq 0}
\end{matrix}\]

写作

\[\mathit{X} \sim \text{Exp}(\mathit{\lambda})\]

\end{definition}

刚才已经计算出, exponential random variable 的 distribution 为

\[\begin{matrix}
{\mathit{F}_{\mathit{X}}(\mathit{x}) = \{\int_{- \infty}^{\mathit{x}}\mathit{f}(\mathit{t})\mathit{d}\mathit{t} = \mathit{\lambda}\int_{0}^{\mathit{x}}\mathit{e}^{- \mathit{\lambda}\mathit{t}}\mathit{d}\mathit{t} = 1 - \mathit{e}^{- \mathit{\lambda}\mathit{x}},\quad\mathit{x} > 0} \\
{0,\quad\mathit{x} < 0}
\end{matrix}\]

容易验证, \(\mathit{F}_{\mathit{X}}(\mathit{x}) = 0\) when \(\mathit{x}\rightarrow - \infty\), 以及 \(\mathit{F}_{\mathit{X}}(\mathit{x})\rightarrow 1\) when \(\mathit{x}\rightarrow\infty\)

计算其 expectation:

\[\begin{matrix}
{\mathbb{E}\lbrack\mathit{X}^{\mathit{n}}\rbrack} & {= \mathit{\lambda}\int_{0}^{\infty}\mathit{x}^{\mathit{n}}\mathit{e}^{- \mathit{\lambda}\mathit{x}}\mathit{d}\mathit{x}} \\
 & {= - \int_{0}^{\infty}\mathit{x}^{\mathit{n}}(\mathit{e}^{- \mathit{\lambda}\mathit{x}})'\mathit{d}\mathit{x}} \\
 & {= 0 + \int_{0}^{\infty}\mathit{n}\mathit{x}^{\mathit{n} - 1}\mathit{e}^{- \mathit{\lambda}\mathit{x}}\mathit{d}\mathit{x}} \\
 & {= \frac{\mathit{n}}{\mathit{\lambda}}\mathbb{E}(\mathit{X}^{\mathit{n} - 1})}
\end{matrix}\]

并 notice: \(\mathbb{E}\lbrack\mathit{X}^{0}\rbrack = \mathbb{E}\lbrack 1\rbrack = 1\). 因而 recursively get:

\[\mathbb{E}\lbrack\mathit{X}^{\mathit{n}}\rbrack = \frac{\mathit{n}!}{\mathit{\lambda}^{\mathit{n}}}\]

即: \(\mathbb{E}\lbrack\mathit{X}\rbrack = \frac{1}{\mathit{\lambda}},\mathbb{E}\lbrack\mathit{X}^{2}\rbrack = \frac{2}{\mathit{\lambda}^{2}},\cdots\) 从而:

\[\mathit{V}\mathit{a}\mathit{r}\lbrack\mathit{X}\rbrack = \frac{2}{\mathit{\lambda}^{2}} - \frac{1}{\mathit{\lambda}^{2}} = \frac{1}{\mathit{\lambda}^{2}}\]

\begin{figure}[htbp]
\centering
\csname prob-02-random-variables-diagram-06\endcsname

\caption{指数分布的概率密度函数}

\label{fig-prob-02-random-variables-diagram-06}

\end{figure}

\begin{figure}[htbp]
\centering
\csname prob-02-random-variables-diagram-07\endcsname

\caption{指数分布的累积分布函数}

\label{fig-prob-02-random-variables-diagram-07}

\end{figure}

\begin{qlremark}{Remark}

exponential 随机变量是 geometric 随机变量的 continuous analogue. 因为 geometric random variable 是假设一个事件在 每个时间点以 constant 的概率发生, 用来 model 它第一次发生的时间;

而 exponential random variable 其实也是, 只是 on a continuous time scale: 它假设的 constant 损坏率其实 implies: \textbf{在任意时间, 如果灯还没坏, 那么灯此刻坏掉的概率是 constant 的.}

\begin{theorem}{\kn{memoryless property of exponential distribution}}

任意 continuous RV,

\[\mathbb{P}(\mathit{X} > \mathit{b} + \mathit{t} \mid \mathit{X} > \mathit{b}) = \mathbb{P}(\mathit{X} > \mathit{t})\Leftrightarrow\mathit{X}\ \text{is an exponential random variable}\]

\end{theorem}

\begin{proof}

简单而言

\[\begin{matrix}
 & {\mathit{P}(\mathit{X} > \mathit{b} + \mathit{t} \mid \mathit{X} > \mathit{b}) = \mathit{P}(\mathit{X} > \mathit{t})} \\
\Leftrightarrow & {\mathit{P}(\mathit{X} > \mathit{b} + \mathit{t}) = \mathit{P}(\mathit{X} > \mathit{t})\mathit{P}(\mathit{X} > \mathit{b})} \\
\Leftrightarrow & {\mathit{X}\ \text{exponential}}
\end{matrix}\]

\end{proof}

因而对于 geometric RV 的 pmf, 当我们把时间间隔 (discrete time interval) 变得越来越小, 这一行为的极限就变为了 exponential RV 的 pdf.

这是一个很自然的 limit behavior, 此处不验证了.

\end{qlremark}

\begin{example}

假设我们有一个 storage battery, 它的 lifetime 是 exponentially distributed 的, 并且 average 为 10 hours. Suppose 我们想要 use this battery for 5 hours for.

计算: probability of finishing the task, if:\\
(a) 使用一个 new battery\\
(b) 使用一个已经用过 2 hours 的 battery\\

\end{example}

\begin{qlsolution}{Solution}

\[\mathbb{E}\lbrack\mathit{X}\rbrack = 10 = \frac{1}{\mathit{\lambda}}\Longrightarrow\mathit{\lambda} = \frac{1}{10}\]

如果我们使用一个 new battery, 则 want: \(\mathit{P}(\mathit{X} > 5) = 1 - \mathit{P}(\mathit{X} < 5) = 1 - \mathit{F}_{\mathit{X}}(5) = \mathit{e}^{- \frac{1}{2}}\) 如果我们使用一个用过 2 hours 的 battery, 则 want:

\[\mathit{P}(\mathit{X} > 5 + 2 \mid \mathit{X} > 2) = \frac{\mathit{e}^{- (5 + 2)\mathit{\lambda}}}{\mathit{e}^{- 2\mathit{\lambda}}} = \mathit{e}^{- 5\mathit{\lambda}} = = \mathit{e}^{- \frac{1}{2}}\]

Notice:

\[\mathit{P}(\mathit{X} > 5) = \mathit{P}(\mathit{X} > 5 + 2 \mid \mathit{X} > 2)\]

\end{qlsolution}

\subsection{Gamma distribution: negative binomial distribution 的 continuous version}\label{gamma-distribution-negative-binomial-distribution-ux7684-continuous-version}

recall Gamma 函数:

\[\Gamma(\mathit{\alpha}) = \int_{0}^{\infty}\mathit{x}^{\mathit{\alpha} - 1}\mathit{e}^{- \mathit{x}}\mathit{d}\mathit{x}\]

这个函数有一个重要的性质:

\[\Gamma(\mathit{\alpha} + 1) = \mathit{\alpha}\Gamma(\mathit{\alpha})\]

因而对于 \(\mathit{n} \in \mathbb{N}\), 我们有:

\[\Gamma(\mathit{n}) = (\mathit{n} - 1)!\]

Gamma 函数就是 factorial 函数的 continuous generalization.

\begin{definition}{\kn{$\Gamma$-distribution}}

如果 continuous random variable \(\mathit{X}\) 的 pdf 是:

\[\mathit{f}_{\mathit{X}}(\mathit{x}) = \left\{ \begin{matrix}
{\frac{\mathit{\lambda}^{\mathit{\alpha}}}{\Gamma(\mathit{\alpha})}\mathit{x}^{\mathit{\alpha} - 1}\mathit{e}^{- \mathit{\lambda}\mathit{x}},} & {\mathit{x} > 0,} \\
{0,} & {\mathit{x} \leq 0}
\end{matrix} \right.\]

那么我们称 \(\mathit{X}\) 服从 \textbf{\(\Gamma\) distribution}, 写作

\[\mathit{X} \sim \Gamma(\mathit{\alpha},\mathit{\lambda})\]

其中 \(\mathit{\alpha} > 0\) 称为 shape parameter, \(\mathit{\lambda} > 0\) 称为 rate parameter.

\end{definition}

\(\Gamma\) distribution 的 cdf 这样求: for \(\mathit{a} \geq 0\),

\[\begin{matrix}
{\mathit{F}_{\mathit{X}}(\mathit{a}) = \int_{- \infty}^{\mathit{a}}\mathit{f}_{\mathit{X}}(\mathit{x})\mathit{d}\mathit{x}} & {= \int_{0}^{\mathit{a}}\frac{\mathit{\lambda}^{\mathit{\alpha}}}{\Gamma(\mathit{\alpha})}\mathit{x}^{\mathit{\alpha} - 1}\mathit{e}^{- \mathit{\lambda}\mathit{x}}\mathit{d}\mathit{x}} \\
 & {= \frac{1}{\Gamma(\mathit{\alpha})}\int_{0}^{\mathit{\lambda}\mathit{a}}\mathit{t}^{\mathit{\alpha} - 1}\mathit{e}^{- \mathit{t}}\mathit{d}\mathit{t}}
\end{matrix}\]

and for \(\mathit{a} < 0\), \(\mathit{F}_{\mathit{X}}(\mathit{a}) = 0\).

Expectation:

\[\begin{matrix}
{\mathbb{E}\lbrack\mathit{X}\rbrack} & {= \frac{1}{\Gamma(\mathit{\alpha})}\int_{0}^{\infty}(\mathit{\lambda}\mathit{x})^{\mathit{\alpha} - 1}\mathit{\lambda}\mathit{e}^{- \mathit{\lambda}\mathit{x}}\mathit{d}\mathit{x}} \\
 & {= \frac{1}{\mathit{\lambda}\Gamma(\mathit{\alpha})}\int_{0}^{\infty}\mathit{x}^{\mathit{\alpha}}\mathit{e}^{- \mathit{x}}\mathit{d}\mathit{x}} \\
 & {= \frac{\Gamma(\mathit{\alpha} + 1)}{\mathit{\lambda}\Gamma(\mathit{\alpha})} = \frac{\mathit{\alpha}}{\mathit{\lambda}}}
\end{matrix}\]

(since \(\Gamma(\mathit{\alpha} + 1) = \mathit{\alpha}\Gamma(\mathit{\alpha})\).) 同理我们可以计算得: \(\mathbb{E}\lbrack\mathit{X}^{2}\rbrack = \frac{\mathit{\alpha}(\mathit{\alpha} + 1)}{\mathit{\lambda}^{2}}\) 从而

\[\text{Var}\lbrack\mathit{X}\rbrack = \mathbb{E}\lbrack\mathit{X}^{2}\rbrack - (\mathbb{E}\lbrack\mathit{X}\rbrack)^{2} = \frac{\mathit{\alpha}(\mathit{\alpha} + 1)}{\mathit{\lambda}^{2}} - \frac{\mathit{\alpha}^{2}}{\mathit{\lambda}^{2}} = \frac{\mathit{\alpha}}{\mathit{\lambda}^{2}}\]

容易发现: 当 \(\mathit{\alpha} = 1\) 时, \(\Gamma\) distribution 就退化成了 exponential distribution.

\[\Gamma(1,\mathit{\lambda}) = \text{Exp}(\mathit{\lambda})\]

而当 \(\mathit{\alpha}\) 是一个正整数 \(\mathit{n}\) 时, \(\Gamma\) distribution 实际就是 \(\mathit{n}\) 个 independent exponential distribution 的 sum:

\begin{proposition}{\kn{Gamma distribution as a sum of independent exponentials}}

Let \(\mathit{X}_{1},\mathit{X}_{2},\cdots,\mathit{X}_{\mathit{n}} \sim \text{Exp}(\mathit{\lambda})\) be i.i.d., 则

\[\Gamma(\mathit{n},\mathit{\lambda}) = \sum\limits_{\mathit{i} = 1}^{\mathit{n}}\mathit{X}_{\mathit{i}}\]

\end{proposition}

\begin{proof}

证明这一结论除了硬算之外, 更快的方法是通过 moment generating function. //TODO

\end{proof}

\begin{qlremark}{Remark}

Gamma distribution 即多个 i.i.d. exponential distribution 的 sum; 而我们知道 exponential distribution 是 geometric distribution 的 continuous analogue; 因而 \textbf{Gamma distribution 也是 negative binomial distribution 的 continuous analogue.}

考虑固定 rate parameter \(\mathit{\lambda} = 1\)，不同 shape parameter \(\mathit{\alpha}\) 的 Gamma distribution：

\begin{figure}[htbp]
\centering
\csname prob-02-random-variables-diagram-08\endcsname

\caption{固定 \(\mathit{\lambda} = 1\) 时不同 shape 参数的 Gamma 密度}

\label{fig-prob-02-random-variables-diagram-08}

\end{figure}

\begin{figure}[htbp]
\centering
\csname prob-02-random-variables-diagram-09\endcsname

\caption{固定 \(\mathit{\lambda} = 1\) 时不同 shape 参数的 Gamma 分布函数}

\label{fig-prob-02-random-variables-diagram-09}

\end{figure}

\begin{itemize}
\item
  当 \(\mathit{\alpha} = 1\) 时, Gamma 分布退化为指数分布, pdf 单调递减;
\item
  随着 \(\mathit{\alpha}\) 增大, pdf 从 0 处开始上升, 形成驼峰形状, 并逐渐向右平移;
\item
  cdf 随 \(\mathit{\alpha}\) 增大而上升变缓, 表示分布的右尾变长。
\end{itemize}

\end{qlremark}

\subsection{normal random variables: 最重要的 continuous RV, 任意 RV 的叠加极限}\label{normal-random-variables-ux6700ux91cdux8981ux7684-continuous-rv-ux4efbux610f-rv-ux7684ux53e0ux52a0ux6781ux9650}

\begin{definition}

我们称 \(\mathit{X}\) is \textbf{normally distributed with parameter \(\mathit{\mu}\) and \(\mathit{\sigma}\)}, if the pdf of \(\mathit{X}\) is given by:

\[\mathit{p}(\mathit{x}) = \frac{1}{\mathit{\sigma}\sqrt{2\mathit{\pi}}}\mathit{e}^{- \frac{(\mathit{x} - \mathit{\mu})^{2}}{2\mathit{\sigma}^{2}}}\]

写作:

\[\mathit{X} \sim \mathcal{N}(\mathit{\mu},\mathit{\sigma}^{2})\]

特别地, 当 \(\mathit{\mu} = 0,\mathit{\sigma} = 1\) 时,

\[\mathit{X} \sim \mathcal{N}(0,1)\]

被称为 \textbf{standard normal distribution.}

\end{definition}

\begin{itemize}
\item
  验证其为 valid pdf:

  我们首先验证 standard normal distribution. 令:

  \[\mathit{I} := \int_{0}^{\infty}\mathit{e}^{\frac{- \mathit{y}^{2}}{2}}\mathit{d}\mathit{y}\]

  则

  \[\begin{matrix}
  \mathit{I}^{2} & {= \int_{0}^{\infty}\mathit{e}^{\frac{- \mathit{y}^{2}}{2}}\mathit{d}\mathit{y} \cdot \int_{0}^{\infty}\mathit{e}^{\frac{- \mathit{x}^{2}}{2}}\mathit{d}\mathit{x}} \\
   & {= \int_{\mathbb{R}_{\geq 0}^{2}}\mathit{e}^{- \frac{\mathit{x}^{2} + \mathit{y}^{2}}{2}}\mathit{d}\mathit{x}\,\mathit{d}\mathit{y}} \\
   & {= \int_{0}^{\mathit{\pi}/2}\int_{0}^{\infty}\mathit{e}^{- \frac{\mathit{r}^{2}}{2}}\mathit{r}\mathit{d}\mathit{r}\,\mathit{d}\mathit{\theta}} \\
   & {= \int_{0}^{\mathit{\pi}/2}( - \mathit{e}^{- \frac{\mathit{r}^{2}}{2}}\ |_{0}^{\infty})\mathit{d}\mathit{\theta}} \\
   & {= \int_{0}^{\mathit{\pi}/2}1\mathit{d}\mathit{\theta} = \frac{\mathit{\pi}}{2}}
  \end{matrix}\]

  从而得到 \(\mathit{I} = \frac{\sqrt{2\mathit{\pi}}}{2}\). 因而 \(\int_{- \infty}^{\infty}\mathit{e}^{- \frac{\mathit{x}^{2}}{2}}\,\mathit{d}\mathit{x} = \sqrt{2\mathit{\pi}}\), 因而 \(\int_{\mathbb{R}}\mathit{p}(\mathit{x})\,\mathit{d}\mathit{x} = 1\).

  从而对于任意的 \(\mathit{\mu}\) 和 \(\mathit{\sigma}\), 通过 change of variable formula 易得 \(\int_{\mathbb{R}}\mathit{p}(\mathit{x})\,\mathit{d}\mathit{x} = 1\).
\item
  计算其 expectation 和 variance:

  我们只需要计算 standard normal distribution 的 expectation 和 variance, 然后通过 expectation 的 linearity 和 variance 的 scale invariance 来计算.

  令 \(\mathit{X} \sim \mathcal{N}(0,1)\). Then

  \[\mathbb{E}\lbrack\mathit{X}\rbrack = \frac{1}{\sqrt{2\mathit{\pi}}}\int_{\mathbb{R}}\mathit{x}\mathit{e}^{- \frac{\mathit{x}^{2}}{2}}\mathit{d}\mathit{x} = 0\]

  因为这是一个 \textbf{odd function}. 从而

  \[\begin{matrix}
  {\mathbb{E}\lbrack\mathit{X}^{2}\rbrack} & {= \frac{1}{\sqrt{2\mathit{\pi}}}\int_{\mathbb{R}}\mathit{x}^{2}\mathit{e}^{- \frac{\mathit{x}^{2}}{2}}\mathit{d}\mathit{x} = \frac{1}{\sqrt{2\mathit{\pi}}}\int_{\mathbb{R}}\mathit{x}(\mathit{x}\mathit{e}^{- \frac{\mathit{x}^{2}}{2}})\mathit{d}\mathit{x} = \frac{1}{\sqrt{2\mathit{\pi}}}\int_{\mathbb{R}}\mathit{x}(\mathit{e}^{- \frac{\mathit{x}^{2}}{2}})'\mathit{d}\mathit{x}} \\
   & {= - \frac{1}{\sqrt{2\mathit{\pi}}}\mathit{x}\mathit{e}^{- \frac{\mathit{x}^{2}}{2}}\ |_{- \infty}^{\infty} + \frac{1}{\sqrt{2\mathit{\pi}}}\int_{\mathbb{R}}\mathit{e}^{- \frac{\mathit{x}^{2}}{2}}\mathit{d}\mathit{x} = 0 + 1 = 1}
  \end{matrix}\]

  从而:

  \[\text{Var}(\mathit{X}) = 1\]

  For general normally distributed \(\mathit{Y}\), 我们知道了 \(\mathit{X}: = \frac{\mathit{Y} - \mathit{\mu}}{\mathit{\sigma}}\) 的 \(\mathbb{E}\lbrack\mathit{X}\rbrack = 0\), \(\text{Var}(\mathit{X}) = 1\), 因而

  \[\mathbb{E}\lbrack\mathit{Y}\rbrack = \mathbb{E}\lbrack\mathit{\sigma}\mathit{X} + \mathit{\mu}\rbrack = \mathit{\mu}\]

  and

  \[\text{Var}(\mathit{Y}) = \text{Var}(\mathit{\sigma}\mathit{X} + \mathit{\mu}) = \mathbb{E}\lbrack(\mathit{\sigma}\mathit{X} + \mathit{\mu} - \mathit{\mu})^{2}\rbrack = \mathbb{E}(\mathit{\sigma}^{2}\mathit{X}^{2}) = \mathit{\sigma}^{2}\mathbb{E}\lbrack\mathit{X}^{2}\rbrack = \mathit{\sigma}^{2}\]

  对于多个 i.i.d. normally distributed \(\mathit{X}\),

  \[\mathbb{E}\lbrack\mathit{X}^{\mathit{k}}\rbrack = \frac{1}{\sqrt{2\mathit{\pi}}}\int_{\mathbb{R}}\mathit{x}^{\mathit{k}}\mathit{e}^{- \frac{\mathit{x}^{2}}{2}}\mathit{d}\mathit{x}\]

  For \(\mathit{k}\) odd, it is \(0\).\\
  For \(\mathit{k}\) even:

  \[\begin{matrix}
  {\mathbb{E}\lbrack\mathit{X}^{\mathit{k}}\rbrack} & {= \frac{1}{\sqrt{2\mathit{\pi}}}\int_{\mathbb{R}}\mathit{x}^{\mathit{k} - 1}(\mathit{x}\mathit{e}^{- \frac{\mathit{x}^{2}}{2}})\mathit{d}\mathit{x}} \\
   & {= - \frac{1}{\sqrt{2\mathit{\pi}}}\int_{\mathbb{R}}\mathit{x}^{\mathit{k} - 1}(\mathit{e}^{- \frac{\mathit{x}^{2}}{2}})'\mathit{d}\mathit{x}} \\
   & {= 0 + \frac{1}{\sqrt{2\mathit{\pi}}}\int_{\mathbb{R}}(\mathit{x}^{\mathit{k} - 1})'\mathit{e}^{- \frac{\mathit{x}^{2}}{2}}\mathit{d}\mathit{x}} \\
   & {= \frac{\mathit{k} - 1}{\sqrt{2\mathit{\pi}}}\int_{\mathbb{R}}\mathit{x}^{\mathit{k} - 2}\mathit{e}^{- \frac{\mathit{x}^{2}}{2}}\mathit{d}\mathit{x}} \\
   & {= (\mathit{k} - 1)\mathbb{E}\lbrack\mathit{X}^{\mathit{k} - 2}\rbrack}
  \end{matrix}\]

  从而:

  \[\begin{matrix}
  {\mathbb{E}\lbrack\mathit{X}^{\mathit{k}}\rbrack = \{ 0,\quad\mathit{k} = 2\mathit{j} - 1} \\
  {(2\mathit{j} - 1)(2\mathit{j} - 3)\cdots 1 = \frac{(2\mathit{j})!}{2^{\mathit{j}}\mathit{j}!},\quad\mathit{k} = 2\mathit{j}}
  \end{matrix}\]
\item
  其 cdf:

  \[\Phi(\mathit{a}) := \mathit{F}_{\mathit{X}}(\mathit{a}) = \int_{( - \infty,\mathit{a}\rbrack}\mathit{\rho}(\mathit{x})\mathit{d}\mathit{x} = \frac{1}{\sqrt{2\mathit{\pi}}}\int_{( - \infty,\mathit{a}\rbrack}\mathit{e}^{- \frac{\mathit{x}^{2}}{2}}\mathit{d}\mathit{x}\]

  我们在微积分中知道, 这个 integral 没有 closed form solution. 因而我们只能给出数值 approximation. Important numerical values:

  \[\Phi( - 3) \approx 0.0013;\quad\Phi( - 2) \approx 0.023;\quad\Phi( - 1) \approx 0.159\]

  Note 由于 \(\mathit{\rho}\) (std) 是一个 even function, 有 \(\mathit{P}(\mathit{X} < - 1) = \mathit{P}(\mathit{X} > 1)\) 因而

  \[\mathit{P}(\mathit{X}\ \text{is one std dev from mean}) = 2(\mathit{P}(\mathit{X} < - 1)) = 2\Phi( - 1) \approx 0.32 = 32\%\]

  Similarly,

  \[\mathit{P}(\mathit{X}\ \text{is two std dev from mean}) \approx 0.046 = 4.6\%\]\[\mathit{P}(\mathit{X}\ \text{is three std dev from mean}) \approx 0.0026 = 0.26\%\]
\end{itemize}

为什么 normal random variables 非常重要: 因为有一个 universality property called \textbf{central limit theorem.} 之后的章节我们会展开讨论:

\paragraph{\texorpdfstring{任意的 random variable, 只要 expectation 和 variance 都是 finite 的, 那么它们的叠加极限都会服从 normal distribution.\\
}{任意的 random variable, 只要 expectation 和 variance 都是 finite 的, 那么它们的叠加极限都会服从 normal distribution. }}\label{ux4efbux610fux7684-random-variable-ux53eaux8981-expectation-ux548c-variance-ux90fdux662f-finite-ux7684-ux90a3ux4e48ux5b83ux4eecux7684ux53e0ux52a0ux6781ux9650ux90fdux4f1aux670dux4ece-normal-distribution.}

这里我们先给出一个 special case of central limit theorem:

\begin{theorem}{\kn{De Moivre--Laplace theorem}}

令 \(\mathit{S}_{\mathit{n}} \sim \text{Bin}(\mathit{n},\mathit{p})\), 则

\[\lim\limits_{\mathit{n}\rightarrow\infty}\mathit{P}(\frac{\mathit{S}_{\mathit{n}} - \mathbb{E}(\mathit{S}_{\mathit{n}})}{\mathit{\sigma}(\mathit{S}_{\mathit{n}})} \in (\mathit{a},\mathit{b})) = \Phi(\mathit{b}) - \Phi(\mathit{a}) = \frac{1}{\sqrt{2\mathit{\pi}}}\int_{(\mathit{a},\mathit{b})}\mathit{e}^{- \frac{\mathit{x}^{2}}{2}}\mathit{d}\mathit{x}\]

\end{theorem}

(Proof: later chapters.)

\begin{example}

Toss a million times 一个 fair coin. Approximate the prob that we get more than \(501000\) heads:

\[\mathbb{E}(\mathit{S}_{1000000}) = \mathit{n}\mathit{p} = 500000\]\[\mathit{\sigma}(\mathit{S}_{1000000}) = \sqrt{\mathit{n}\mathit{p}\mathit{q}} = 500\]

因而:

\[\begin{matrix}
{\mathbb{P}(\mathit{S}_{1000000} > 501000)} & {= \mathbb{P}(\mathit{S}_{1000000} - \mathbb{E}(\mathit{S}_{1000000}) > 1000)} \\
 & {= \mathbb{P}\left( {\frac{\mathit{S}_{1000000} - \mathbb{E}(\mathit{S}_{1000000})}{\mathit{\sigma}(\mathit{S}_{1000000})} > \frac{1000}{500}} \right)} \\
 & {\approx 1 - \Phi(2) = \Phi( - 2) \approx 0.159}
\end{matrix}\]

\end{example}

\paragraph{\texorpdfstring{除此之外 need to mention: Poission distribution 当 \(\mathit{\lambda}\rightarrow\infty\) 时, 也得到 normal distribution.\\
}{除此之外 need to mention: Poission distribution 当 \textbackslash mathit\{\textbackslash lambda\}\textbackslash rightarrow\textbackslash infty 时, 也得到 normal distribution. }}\label{ux9664ux6b64ux4e4bux5916-need-to-mention-poission-distribution-ux5f53-mathitlambdarightarrowinfty-ux65f6-ux4e5fux5f97ux5230-normal-distribution.}

这并不是一个令人惊讶的事情. 也可以用 central limit theorem 来解释. 因为 recall: Poisson distribution 有一个很重要的 property: additivity. 因而 \(\mathit{\lambda}\rightarrow\infty\) 也可以看作 \(\mathit{n}\) 个 i.i.d. Poisson distribution 的 sum 在 \(\mathit{n}\rightarrow\infty\) 时的极限行为, 从而通过 central limit theorem 得到 normal distribution.

\begin{landscape}
\subsection{summmary of discrete and continuous variables examples}\label{summmary-of-discrete-and-continuous-variables-examples}

{\def\LTcaptype{none} % do not increment counter
\footnotesize\setlength{\tabcolsep}{3pt}\renewcommand{\arraystretch}{1.4}
\begin{longtable}[]{@{}p{0.12\linewidth}p{0.25\linewidth}p{0.25\linewidth}p{0.06\linewidth}p{0.07\linewidth}p{0.18\linewidth}@{}}
\toprule\noalign{}
\textbf{distribution} & \textbf{pmf / pdf} & \textbf{cdf} & \textbf{\(\mathbb{E}\lbrack\mathit{X}\rbrack\)} & \textbf{\(\text{Var}(\mathit{X})\)} & \textbf{meaning} \\
\midrule\noalign{}
\endhead
\bottomrule\noalign{}
\endlastfoot
\(\text{Ber}(\mathit{p})\) & \(\mathbb{P}(\mathit{X} = 1) = \mathit{p}\), \(\mathbb{P}(\mathit{X} = 0) = 1 - \mathit{p}\) & \(\mathit{F}(\mathit{x}) = 1 - \mathit{p},0 \leq \mathit{x} < 1\) & \(\mathit{p}\) & \(\mathit{p}(1 - \mathit{p})\) & 一次 Bernoulli trial (unfair coin) 的结果 \\
\(\text{Bin}(\mathit{n},\mathit{p})\) & \(\mathbb{P}(\mathit{X} = \mathit{k}) = \binom{\mathit{n}}{\mathit{k}}\mathit{p}^{\mathit{k}}(1 - \mathit{p})^{\mathit{n} - \mathit{k}}\), \(\mathit{k} = 0,\ldots,\mathit{n}\) & \(\mathit{F}(\mathit{k}) = \sum_{\mathit{j} = 0}^{\lfloor\mathit{k}\rfloor}\binom{\mathit{n}}{\mathit{j}}\mathit{p}^{\mathit{j}}(1 - \mathit{p})^{\mathit{n} - \mathit{j}}\) & \(\mathit{n}\mathit{p}\) & \(\mathit{n}\mathit{p}(1 - \mathit{p})\) & \(\mathit{n}\) 次 i.i.d. Bernoulli trials 中的成功次数 \\
\(\text{Geom}(\mathit{p})\) & \(\mathbb{P}(\mathit{X} = \mathit{k}) = (1 - \mathit{p})^{\mathit{k} - 1}\mathit{p}\), \(\mathit{k} \geq 1\) & \(\mathit{F}(\mathit{k}) = 1 - (1 - \mathit{p})^{\mathit{k}}\) & \(\frac{1}{\mathit{p}}\) & \(\frac{1 - \mathit{p}}{\mathit{p}^{2}}\) & Bernoulli trial 第一次成功所需的试验次数 \\
\(\text{NB}(\mathit{r},\mathit{p})\) & \(\mathbb{P}(\mathit{X} = \mathit{k}) = \binom{\mathit{k} + \mathit{r} - 1}{\mathit{k}}(1 - \mathit{p})^{\mathit{k}}\mathit{p}^{\mathit{r}},\ \mathit{k} \geq 0\) & \(\mathit{F}(\mathit{k}) = \sum_{\mathit{j} = 0}^{\mathit{k}}\binom{\mathit{j} + \mathit{r} - 1}{\mathit{j}}(1 - \mathit{p})^{\mathit{j}}\mathit{p}^{\mathit{r}}\) & \(\frac{\mathit{r}(1 - \mathit{p})}{\mathit{p}}\) & \(\frac{\mathit{r}(1 - \mathit{p})}{\mathit{p}^{2}}\) & Bernoulli trial 在 \(\mathit{r}\) 次成功之前观察到的失败次数 \\
\(\text{Poi}(\mathit{\lambda})\) & \(\mathbb{P}(\mathit{X} = \mathit{k}) = \mathit{e}^{- \mathit{\lambda}}\frac{\mathit{\lambda}^{\mathit{k}}}{\mathit{k}!},\ \mathit{k} \geq 0\) & \(\mathit{F}(\mathit{k}) = \sum_{\mathit{j} = 0}^{\mathit{k}}\mathit{e}^{- \mathit{\lambda}}\frac{\mathit{\lambda}^{\mathit{j}}}{\mathit{j}!}\) & \(\mathit{\lambda}\) & \(\mathit{\lambda}\) & fixed window (time/space) 下某一事件的计数 \\
\(\mathit{U}(\lbrack\mathit{a},\mathit{b}\rbrack)\) & \(\mathit{f}(\mathit{x}) = \frac{1}{\mathit{b} - \mathit{a}}\), \(\mathit{a} \leq \mathit{x} \leq \mathit{b}\) & \(\mathit{F}(\mathit{x}) = \left\{ \begin{matrix}
{0,} & {\mathit{x} < \mathit{a}} \\
{\frac{\mathit{x} - \mathit{a}}{\mathit{b} - \mathit{a}},} & {\mathit{a} \leq \mathit{x} \leq \mathit{b}} \\
{1,} & {\mathit{x} > \mathit{b}}
\end{matrix} \right.\) & \(\frac{\mathit{a} + \mathit{b}}{2}\) & \(\frac{(\mathit{b} - \mathit{a})^{2}}{12}\) & uniform choice on an interval \\
\(\text{Exp}(\mathit{\lambda})\) & \(\mathit{f}(\mathit{x}) = \mathit{\lambda}\mathit{e}^{- \mathit{\lambda}\mathit{x}}\), \(\mathit{x} \geq 0\) & \(\mathit{F}(\mathit{x}) = \left\{ \begin{matrix}
{0,} & {\mathit{x} \leq 0} \\
{1 - \mathit{e}^{- \mathit{\lambda}\mathit{x}},} & {\mathit{x} > 0}
\end{matrix} \right.\) & \(\frac{1}{\mathit{\lambda}}\) & \(\frac{1}{\mathit{\lambda}^{2}}\) & 固定平均发生率下, 某事件第一次发生要等多久 \\
\(\Gamma(\mathit{\alpha},\mathit{\lambda})\) & \(\mathit{f}(\mathit{x}) = \frac{\mathit{\lambda}^{\mathit{\alpha}}}{\Gamma(\mathit{\alpha})}\mathit{x}^{\mathit{\alpha} - 1}\mathit{e}^{- \mathit{\lambda}\mathit{x}}\), \(\mathit{x} \geq 0\) & \(\mathit{F}(\mathit{x}) = \frac{1}{\Gamma(\mathit{\alpha})}\int_{0}^{\mathit{\lambda}\mathit{x}}\mathit{t}^{\mathit{\alpha} - 1}\mathit{e}^{- \mathit{t}}\,\mathit{d}\mathit{t}\) & \(\frac{\mathit{\alpha}}{\mathit{\lambda}}\) & \(\frac{\mathit{\alpha}}{\mathit{\lambda}^{2}}\) & \(\mathit{\alpha}\) 个 i.i.d. exponential distributions 的和 (固定平均发生率下，第 \(\mathit{\alpha}\) 次事件出现前要等多久) \\
\(\mathit{N}(\mathit{\mu},\mathit{\sigma}^{2})\) & \(\mathit{f}(\mathit{x}) = 1/\mathit{\sigma}\sqrt{2\mathit{\pi}}\exp( - (\mathit{x} - \mathit{\mu})^{2}/2\mathit{\sigma}^{2})\) & no elementary closed form & \(\mathit{\mu}\) & \(\mathit{\sigma}^{2}\) & 大量微小独立扰动 (测量误差/自然波动) 叠加后的结果 \\
\end{longtable}
}
\end{landscape}
